% =====================================================================
% 第 6 章 Wakimoto 模
% 覆盖：PDF p177--p208（原书印刷 p.165--p.196）
% 底本：Frenkel, Langlands Correspondence for Loop Groups, CUP 2007
% =====================================================================
%
% 编号说明：preamble.tex 中定理类环境分两组按节编号（定理与引理共用、
% 命题与推论共用），而原书第 6 章的所有定理类陈述共用同一个按节计数器
% （如 引理 6.1.1、推论 6.1.2、定理 6.1.3、…… 连续编号）。自动编号无法
% 复现原书编号，故下面在每条陈述之前按原印刷编号逐一设定该环境的编号
% 显示。整章包在 \begingroup...\endgroup 中，文件末尾自动恢复 preamble
% 的定义。
%
\begingroup

\chapter{Wakimoto 模}

本章把前面构造的自由场实现推广到非临界水平，并用顶点代数的语言重述它。
用这些术语，自由场实现就是从与 $\gghat_\kappa$ 相伴的顶点代数
$V_\kappa(\fg)$ 到与无穷维 Heisenberg 李代数相伴的“自由场”顶点代数
$M_{\mathfrak g}\otimes\pi_0^{\kappa-\kappa_c}$ 的一个同态。

顶点代数 $M_{\mathfrak g}\otimes\pi_0^{\kappa-\kappa_c}$ 上的模现在是
$\gghat_\kappa$-模，我们称之为 \textbf{Wakimoto 模}。它们最初由 M. Wakimoto
[1986] 在 $\fg=\widehat{\fsl}_2$ 情形构造，并由 B. Feigin 与作者
[1988; 1990a] 对一般 $\fg$ 构造。Wakimoto 模可以看作 $\gghat$ 的表示，
它们是从其 Heisenberg 子代数 $\widehat{\mathfrak h}$ 的表示“半无限诱导”得到的。
然而与通常的诱导不同，$\widehat{\mathfrak h}$-模的水平被临界值 $\kappa_c$ 平移。
特别地，若从水平为零的 $\widehat{\mathfrak h}$-模（例如对应阿贝尔李代数
$L\mathfrak h$ 的一个特征标的一维模）出发，所得的 $\gghat$-模就在临界水平。
这就是为什么“中点”不是零水平（如人们可能天真预期的那样），而是临界水平。

Wakimoto 模（临界与非临界水平的）对实现最终目标至关重要：描述
$V_{\kappa_c}(\fg)$ 的中心（见第 8 章）。

本章内容更详细的描述如下。我们在 \S6.1 从临界水平情形开始。主要结果是
顶点代数同态 $V_{\kappa_c}(\fg)\to M_{\mathfrak g}\otimes\pi_0$ 的存在性，
其中 $\pi_0$ 是与 $L\mathfrak h$ 相伴的交换顶点代数。利用这个同态，我们在
张量积 $M_{\mathfrak g}\otimes N$ 上构造临界水平的 $\gghat$-模结构，其中
$N$ 是任意（光滑）$L\mathfrak h$-模。这些就是临界水平的 Wakimoto 模（下文
将看到，对参数的一般取值它们不可约）。

在 \S6.2 中把构造推广到任意水平。我们证明顶点代数同态
\[
    w_\kappa:V_\kappa(\fg)\to M_{\mathfrak g}\otimes\pi_0^{\kappa-\kappa_c}
\]
的存在性，其中 $\pi_0^{\kappa-\kappa_c}$ 是与 Heisenberg 李代数
$\widehat{\mathfrak h}_{\kappa-\kappa_c}$ 相伴的顶点代数。与前面讨论中心时一样，
我们特别注意坐标变换群 $\operatorname{Aut}\mathcal O$ 在 Wakimoto 模上的
作用。这是为了得到这些模的无坐标构造，而中心的描述正需要无坐标描述。为此，
我们描述 $M_{\mathfrak g}\otimes\pi_0^{\kappa-\kappa_c}$ 上与
$V_\kappa(\fg)$ 上的 Segal--Sugawara（拟）共形结构相对应的（拟）共形结构。

接着在 \S6.3 中按 [Feigin and Frenkel 1990a] 的思想，把 Wakimoto 模的
构造推广到更一般的背景：把李子代数 $\widehat{\mathfrak h}$ 换成任意抛物子代数
$\mathfrak p$ 的 Levi 子代数的圈代数的中心扩张。于是我们建立了仿射
Kac--Moody 代数表示的“半无限抛物诱导”模式，类似于局部非阿基米德域上约化群的
抛物诱导。

\begin{keypoints}
    \item \textbf{核心思想：} Wakimoto 模把 $\gghat_\kappa$-模实现为自由场
    （Weyl 代数 + Heisenberg 代数）顶点代数上的模，其水平相对 $\kappa_c$ 有
    平移；这是“半无限诱导”。
    \item \textbf{数学工具：} 顶点代数同态、正规序、Segal--Sugawara 向量的像、
    拟共形结构、$\operatorname{Aut}\mathcal O$-torsor、Chevalley 复形与上同调。
    \item \textbf{关键结论：} 存在 $V_{\kappa_c}(\fg)\to M_{\mathfrak g}\otimes
    \pi_0$ 及 $V_\kappa(\fg)\to M_{\mathfrak g}\otimes\pi_0^{\kappa-\kappa_c}$
    的同态；Wakimoto 模可刻画为半无限抛物诱导；Kac--Kazhdan 猜想由此得证。
    \item \textbf{记号提醒：}$a_{\alpha,n}$ 与 $a^*_{\alpha,n}$ 是 Weyl 代数生成元，
    $b_{i,n}$ 是 Heisenberg 代数生成元；$M_{\mathfrak g}$ 与 $\pi_0$ 都是顶点
    代数，$\otimes$ 为顶点代数张量积。
\end{keypoints}

\section{临界水平的 Wakimoto 模}

\setcounter{theorem}{0}

\begin{keypoints}
    \item \textbf{核心思想：} 临界水平的自由场实现是顶点代数同态
    $V_{\kappa_c}(\fg)\to M_{\mathfrak g}\otimes\pi_0$；模同态把
    $\pi_0$-模变成临界水平的 $\gghat$-模。
    \item \textbf{数学工具：} 由生成向量的 Fourier 系数确定的同态（引理 6.1.1）、
    $H^1(L\fg,\cA^{\mathfrak g}_{0,\mathrm{loc}})$ 的权 0 部分、正规序与双缩并。
    \item \textbf{关键结论：} 对每个 $\chi(t)\in\mathfrak h^*(\!(t)\!)$ 得到临界
    水平的 Wakimoto 模 $W_{\chi(t)}$；更强的形式（定理 6.1.6）给出
    $V_{\kappa_c}(\fg)\to M_{\mathfrak g}\otimes\pi_0$ 的显式公式。
    \item \textbf{记号提醒：}$\cA^{\mathfrak g}_{1,\mathrm{loc}}$ 是局部 Weyl 代数，
    $\cA^{\mathfrak g}_{0,\mathrm{loc}}$ 是其常数（权 0）部分；$a^*_{\alpha,n}$ 带
    星号。
\end{keypoints}

本节把上一章得到的自由场实现重述为顶点代数同态
$V_{\kappa_c}(\fg)\to M_{\mathfrak g}\otimes\pi_0$，其中 $\pi_0$ 是与
$L\mathfrak h$ 相伴的交换顶点代数。

\subsection{顶点代数同态}

回忆与仿射 Kac--Moody 代数 $\gghat$ 相伴的真空 $\gghat$-模
$V_\kappa(\fg)$，其定义见 \S2.2.4。定理 2.2.2 已证明 $V_\kappa(\fg)$ 是
$\Z$-分次顶点代数。我们希望把李代数同态
$\gghat_{\kappa_c}\to\cA^{\mathfrak g}_{1,\mathrm{loc}}$ 解释为顶点代数同态。

\renewcommand{\thelemma}{6.1.1}
\begin{lemma}
    定义 $\Z$-分次顶点代数同态 $V_\kappa(\fg)\to V$，等价于选取次数为 $1$ 的
    向量 $\widetilde J^a_{-1}|0\rangle_V\in V$（$a=1,\dots,\dim\fg$），使得顶点
    算子
    \[
        Y(\widetilde J^a_{-1}|0\rangle_V,z)
        =\sum_{n\in\Z}\widetilde J^a_nz^{-n-1}
    \]
    的 Fourier 系数 $\widetilde J^a_n$ 满足关系 (1.3-4)，其中 $\mathbf 1=1$。
\end{lemma}

\begin{proof}
    给定同态 $\rho:V_\kappa(\fg)\to V$，令
    $\widetilde J^a_{-1}|0\rangle_V=\rho(J^a_{-1}|0\rangle)$。$\rho$ 是顶点代数
    同态这一事实说明，顶点算子 $Y(\widetilde J^a_{-1}|0\rangle_V,z)$ 的 OPE 与
    顶点算子 $J^a(z)$ 的相同，从而它们的 Fourier 系数 $\widetilde J^a_n$ 的
    交换关系与 $J^a_n$ 的相同。

    反之，设给定满足引理条件的向量 $\widetilde J^a_{-1}|0\rangle_V$。用公式
    \[
        J^{a_1}_{n_1}\cdots J^{a_m}_{n_m}|0\rangle
        \mapsto\widetilde J^{a_1}_{n_1}\cdots\widetilde J^{a_m}_{n_m}|0\rangle
    \]
    定义线性映射 $\rho:V_\kappa(\fg)\to V$。容易验证它是 $\Z$-分次顶点代数
    同态。
\end{proof}

顶点代数 $M_{\mathfrak g}$ 是 $\Z$-分次的（见 \S5.4.1）。下述结果是
定理 5.6.8 的推论。

\renewcommand{\thecorollary}{6.1.2}
\begin{corollary}
    存在 $\Z$-分次顶点代数同态 $V_{\kappa_c}(\fg)\to M_{\mathfrak g}$。
\end{corollary}

\begin{proof}
    由定理 5.6.8，存在李代数同态 $\gghat_{\kappa_c}\to U(M_{\mathfrak g})$，
    使 $\mathbf 1\mapsto1$ 且
    \[
        J^a(z)\mapsto Y(P_a(a^*_{\alpha,0},a_{\beta,-1})|0\rangle,z)+Y(B_a,z),
    \]
    其中向量 $P_a$ 与 $B_a$ 的次数为 $1$。因此由引理 6.1.1 得到顶点代数同态
    $V_{\kappa_c}(\fg)\to M_{\mathfrak g}$，使得
    \[
        J^a_{-1}|0\rangle\mapsto
        P_a(a^*_{\alpha,0},a_{\beta,-1})|0\rangle+B_a.
    \]
\end{proof}

为证明定理 5.6.8，注意复形 $C^\bullet_{\mathrm{loc}}(L\fg,
\cA^{\mathfrak g}_{0,\mathrm{loc}})$ 带有关于 $\fg$ 的根格的分次，它由公式
\[
    \mathrm{wt}\,a^*_{\alpha,n}=-\alpha,\qquad
    \mathrm{wt}\,\psi^*_{a,n}=-\mathrm{wt}\,J^a,\qquad
    \mathrm{wt}\,|0\rangle=0
\]
定义，使得 $\mathrm{wt}\int Y(A,z)dz=\mathrm{wt}A$。微分保持这个分次，并且
显然上循环 $\omega$ 与 $i_*(\sigma)$ 关于它是齐次的。因此定理 5.6.8 证明中
引入的元素 $\gamma$ 可以选成关于权重分次也是齐次的。这意味着 $B_a\in
M_{\mathfrak g,+}$ 可以选成除 $\deg B_a=1$ 之外还满足
$\mathrm{wt}\,B_a=\mathrm{wt}\,J^a$。\transnote{卡点：此处同时用到两种分次——根格权 $\mathrm{wt}$ 与顶点代数次数 $\deg$，读者容易混为一谈。背景：$\mathrm{wt}\,a^*_{\alpha,n}=-\alpha$、$\mathrm{wt}\,a_{\alpha,n}=\alpha$、$\mathrm{wt}\,b_{i,n}=0$，而 $\deg$ 由 Virasoro 算子 $L_0$ 的本征值给出，两者彼此独立；把 $\gamma$ 取成对权分次齐次之后，$B_a$ 必须同时满足 $\deg B_a=1$ 与 $\mathrm{wt}\,B_a=\mathrm{wt}\,J^a$。落点：\S6.1.1、定理 6.1.3 之前、原书 p.167。来源：原书 p.167 的分次论证；并见 \S6.2.3 证明中两种分次的并用。}

于是对一切 $J^a\in\mathfrak h\oplus\mathfrak n_+$ 必有 $B_a=0$，因为
$M_{\mathfrak g,+}$ 中没有这些权重的元素。进一步，对应于 $J^a=f_i$ 的项
$B_a$ 必与 $a^*_{\alpha_i,-1}|0\rangle\in M_{\mathfrak g,+}$ 成比例。利用
\S5.2.5 的公式与 \S5.3.2 的讨论，便得到同态
$V_{\kappa_c}(\fg)\to M_{\mathfrak g}$ 的更显式描述：
% [Tutor Note: 这里的论证用到“权 0 上循环”与根的权重匹配：$a^*_{\alpha,n}$ 的权为 $-\alpha$，故 $B_a$ 的权必须等于 $\mathrm{wt}\,J^a$；对 $J^a\in\mathfrak h\oplus\mathfrak n_+$ 没有这样权的向量，对 $J^a=f_i$ 唯一可能是 $a^*_{\alpha_i,-1}$。这是第 5 章 Chevalley 复形计算的直接推论，细节见 [Frenkel and Ben-Zvi 2004] 及 [Feigin and Frenkel 1990a]。]

\renewcommand{\thetheorem}{6.1.3}
\begin{theorem}
    存在常数 $c_i\in\C$，使得顶点算子
    \[
        \begin{aligned}
            e_i(z)&=a_{\alpha_i}(z)+\sum_{\beta\in\Delta_+}
            :\!P^i_\beta(a^*_\alpha(z))a_\beta(z)\!:,\\
            h_i(z)&=-\sum_{\beta\in\Delta_+}
            \beta(h_i):\!a^*_\beta(z)a_\beta(z)\!:,\\
            f_i(z)&=\sum_{\beta\in\Delta_+}
            :\!Q^i_\beta(a^*_\alpha(z))a_\beta(z)\!:
            +c_i\partial_za^*_{\alpha_i}(z)
        \end{aligned}
    \]
    的 Fourier 系数生成 $\gghat_{\kappa_c}$ 在 $M_{\mathfrak g}$ 上的作用，
    其中多项式 $P^i_\beta$、$Q^i_\beta$ 由公式 (5.2-12)--(5.2-14) 引入。
\end{theorem}

除上述李代数同态 $w_{\kappa_c}:\gghat_{\kappa_c}\to
\cA^{\mathfrak g}_{1,\mathrm{loc}}$ 外，还有李代数反同态
\[
    w^R:L\mathfrak n_+\to\cA^{\mathfrak g}_{1,\mathrm{loc}},
\]
它由 $\mathfrak n_+$ 在 $N_+$ 上的右作用诱导（见 \S5.2.5）。按构造，
$L\mathfrak n_+$ 在 $w_{\kappa_c}$ 与 $w^R$ 下的像彼此交换。我们有
\begin{equation}
    w^R(e_i(z))=a_{\alpha_i}(z)+\sum_{\beta\in\Delta_+}
    P^{R,i}_\beta(a^*_\gamma(z))a_\beta(z),
    \tag{6.1-1}
\end{equation}
其中多项式 $P^{R,i}_\beta$ 已在 \S5.2.5 中定义。更一般地，有
\begin{equation}
    w_{\kappa_c}(e_\alpha(z))=a_\alpha(z)+\sum_{\beta\in\Delta_+;\beta>\alpha}
    P^\alpha_\beta(a^*_\gamma(z))a_\beta(z),
    \tag{6.1-2}
\end{equation}
\begin{equation}
    w^R(e_\alpha(z))=a_\alpha(z)+\sum_{\beta\in\Delta_+;\beta>\alpha}
    P^{R,\alpha}_\beta(a^*_\gamma(z))a_\beta(z),
    \tag{6.1-3}
\end{equation}
对某些多项式 $P^\alpha_\beta$ 与 $P^{R,\alpha}_\beta$，其中条件
$\beta>\alpha$ 来自交换关系
\[
    [h,e_\alpha]=\alpha(h)e_\alpha,\qquad
    [h,e^R_\alpha]=\alpha(h)e^R_\alpha,\qquad h\in\mathfrak h.
\]
注意上述三个公式中不需要正规序，理由相同：$P^\alpha_\beta$ 与
$P^{R,\alpha}_\beta$ 不可能含 $a_\beta$。\transnote{卡点：(6.1-1)--(6.1-3) 的多项式写作 $a^*_\gamma(z)$，定理 6.1.3 中同类多项式却写作 $a^*_\alpha(z)$，是笔误还是两种记号？背景：这些多项式其实是全体 $a^*_\beta$（$\beta\in\Delta_+$）的多项式，下标 $\alpha$、$\gamma$ 只是"任取一个根"的占位记号，并不指某个特定因子；原书 p.168 与 p.170 分别如此排印，两种写法并存。又因 $P^\alpha_\beta$、$P^{R,\alpha}_\beta$ 不含 $a_\beta$，与 $a_\beta(z)$ 相乘时没有需要正规序的缩并。落点：\S6.1.1 式 (6.1-1)--(6.1-3) 与定理 6.1.3。来源：原书 p.168、p.170；并见本译本 \S5.2.5。}
% [Tutor Note: 原书这三个公式里的多项式参数印作 $a^*_\gamma(z)$（(6.1-1)--(6.1-3)），而定理 6.1.3 中同类多项式印作 $a^*_\alpha(z)$。两种写法在原书中并存，疑为同一“全体 $a^*$ 生成元”的记号简写；此处按各页原文分别照录，未作统一。]

\subsection{$M_{\mathfrak g}$ 上的其他 $\gghat$-模结构}

在定理 6.1.3 中我们构造了 $M_{\mathfrak g}$ 上的 $\gghat_{\kappa_c}$-模结构。
要得到 $M_{\mathfrak g}$ 上其他临界水平的 $\gghat_{\kappa_c}$-模结构，需要
考虑提升同态 $L\fg\to T\fg_{\mathrm{loc}}$ 的其他同态
$\gghat_{\kappa_c}\to\cA^{\mathfrak g}_{1,\mathrm{loc}}$。由引理 5.5.3，这类
提升的同构类构成 $H^1(L\fg,\cA^{\mathfrak g}_{0,\mathrm{loc}})$ 上的一个
torsor，后者是“大”Chevalley 复形
\[
    C^i(L\fg,\cA^{\mathfrak g}_{0,\mathrm{loc}})
    =\Hom_{\mathrm{cont}}\left(\bigwedge\nolimits^i(L\fg),
    \cA^{\mathfrak g}_{0,\mathrm{loc}}\right)
\]
的一阶上同调。回忆我们的复形
$C^\bullet_{\mathrm{loc}}(L\fg,\cA^{\mathfrak g}_{0,\mathrm{loc}})$ 有权的
分次，且我们的上循环 $\omega$ 的权为 $0$。因此在所有提升中，我们考虑权为
$0$ 的那些。这类提升的集合与 $H^1(L\fg,\cA^{\mathfrak g}_{0,\mathrm{loc}})$
的权 $0$ 齐次分量双射。\transnote{卡点：为何"权 0 的提升"恰好对应 $H^1$ 的权 0 分量，而这个分量又被 $(L\mathfrak h)^*$ 刻画？背景：提升之间的差是以 $\cA^{\mathfrak g}_{0,\mathrm{loc}}$ 为系数的一阶上循环，模去上边缘即得 $H^1$；要求权为 $0$ 是为了让提升保持根格分次。引理 6.1.4 进一步断言该分量同构于 $(L\mathfrak h)^*$，于是在临界水平上，$M_{\mathfrak g}$ 上各模结构之间的差别恰好由 $L\mathfrak h$ 上的线性泛函（即 $\mathfrak h$ 流的平移）参数化。落点：\S6.1.1、引理 6.1.4、\S6.1.2。来源：原书 p.168--169（引理 6.1.4 及其证明）；并见 \S5.5 的提升理论。}
% [Tutor Note: 这里是第 5 章“提升”理论的直接应用：$w_{\kappa_c}$ 固定后，所有提升的差别是 $L\fg$ 的一个以 $\cA^{\mathfrak g}_{0,\mathrm{loc}}$ 为系数的一阶上循环，模去上边缘即得 $H^1$。权 0 条件保证提升保持权分次，这正是下文定理 6.1.3 中 $B_a$ 的权必须等于 $\mathrm{wt}\,J^a$ 的来源。]

\renewcommand{\thelemma}{6.1.4}
\begin{lemma}
    $H^1(L\fg,\cA^{\mathfrak g}_{0,\mathrm{loc}})$ 的权 $0$ 分量同构于
    （拓扑）对偶空间 $(L\mathfrak h)^*$。
\end{lemma}

\begin{proof}
    先证明 $H^1(L\fg,\cA^{\mathfrak g}_{0,\mathrm{loc}})$ 的权 $0$ 分量同构于
    $C^1(L\fg,\cA^{\mathfrak g}_{0,\mathrm{loc}})$ 中权 $0$ 上循环构成的空间。
    事实上，由于 $\mathrm{wt}\,a^*_{\alpha,n}=-\alpha$，
    $C^0(L\fg,\cA^{\mathfrak g}_{0,\mathrm{loc}})=\cA^{\mathfrak g}_{0,\mathrm{loc}}$
    的权 $0$ 部分是一维的，由常数组成。对常数作用微分得到 $0$，所以
    $C^1(L\fg,\cA^{\mathfrak g}_{0,\mathrm{loc}})$ 中没有权 $0$ 的上边缘。

    其次证明任意权 $0$ 的一阶上循环 $\phi$ 由它在 $L\mathfrak h\subset L\fg$
    上的限制唯一确定，而这个限制可以是 $L\mathfrak h$ 上任意的（连续）线性
    泛函。事实上，由于 $\cA^{\mathfrak g}_{0,\mathrm{loc}}$ 中出现的权都
    $\le0$，$\phi$ 在 $L_+\mathfrak n_-$ 上的限制等于 $0$，在 $L\mathfrak h$
    上的限制取值于常数 $\C\subset\cA^{\mathfrak g}_{0,\mathrm{loc}}$。现在固定
    $\phi|_{L\mathfrak h}$。用留数配对把 $(L\mathfrak h)^*$ 与
    $\mathfrak h(\!(t)\!)dt$ 等同，并在此等同下把 $\phi|_{L\mathfrak h}$ 写成
    $\chi(t)dt$。这里
    \[
        \chi(t)=\sum_{n\in\Z}\chi_nt^{-n-1},\qquad\chi_n\in\mathfrak h^*,
    \]
    其中 $\chi_n(h)=\phi(h_n)$。把 $\langle\chi(t),h_i\rangle$ 记为
    $\chi_i(t)$。

    我们断言：对任意 $\chi(t)dt\in\mathfrak h^*(\!(t)\!)dt$，存在唯一的权
    $0$ 一阶上循环 $\phi\in C^1(L\fg,\cA^{\mathfrak g}_{0,\mathrm{loc}})$，
    使得
    \begin{equation}
        \phi(e_i(z))=0,
        \tag{6.1-4}
    \end{equation}
    \begin{equation}
        \phi(h_i(z))=\chi_i(z),
        \tag{6.1-5}
    \end{equation}
    \begin{equation}
        \phi(f_i(z))=\chi_i(z)a^*_{\alpha,n}(z).
        \tag{6.1-6}
    \end{equation}
    事实上，固定 (6.1-4) 与 (6.1-5) 中的 $\phi(e_i(z))$、$\phi(h_i(z))$ 之后，
    利用公式
    \[
        \phi([e_{i,n},f_{j,m}])=e_{i,n}\cdot\phi(f_{j,m})
        -f_{j,m}\cdot\phi(e_{i,n})
    \]
    得到
    \[
        \delta_{i,j}\chi_{i,n+m}=e_{i,n}\cdot\phi(f_{j,m}).
    \]
% [Tutor Note: (6.1-6) 与 (6.1-9) 中末项印作 $\chi_i(z)a^*_{\alpha,n}(z)$，而下文 (6.1-10) 对应项印作 $b_i(z)a^*_{\alpha_i}(z)$。按语境应为 $a^*_{\alpha_i}(z)$（$a^*$ 的场），此处按原文照录，并在报告中登记为疑似笔误。]
\end{proof}

关于 $(\phi(f_{j,m}))$ 的这个方程在权 $-\alpha_i$ 的
$\cA^{\mathfrak g}_{0,\mathrm{loc}}$ 中有唯一解，即公式 (6.1-6) 给出的解。
若上循环 $\phi$ 存在，则一旦固定它在 $e_{i,n}$、$h_{i,n}$、$f_{i,n}$ 上的
值就唯一确定。下面证明它存在。这等价于证明场
\begin{equation}
    e_i(z)=a_{\alpha_i}(z)+\sum_{\beta\in\Delta_+}
    :\!P^i_\beta(a^*_\alpha(z))a_\beta(z)\!:,
    \tag{6.1-7}
\end{equation}
\begin{equation}
    h_i(z)=-\sum_{\beta\in\Delta_+}
    \beta(h_i):\!a^*_\beta(z)a^*_\beta(z)\!:+\chi_i(z),
    \tag{6.1-8}
\end{equation}
\begin{equation}
    f_i(z)=\sum_{\beta\in\Delta_+}
    :\!Q^i_\beta(a^*_\alpha(z))a_\beta(z)\!:
    +c_i\partial_za^*_{\alpha_i}(z)+\chi_i(z)a^*_{\alpha,n}(z)
    \tag{6.1-9}
\end{equation}
的 Fourier 系数满足 $\gghat_{\kappa_c}$ 的关系，其中 $\mathbf 1=1$。\transnote{卡点：(6.1-6)、(6.1-9) 的末项印作 $\chi_i(z)a^*_{\alpha,n}(z)$，(6.1-8) 印作 $:a^*_\beta(z)a^*_\beta(z):$，与定理 6.1.3 中 $h_i$ 的 $:a^*_\beta(z)a_\beta(z):$ 不一致，是原书笔误吗？背景：按权与正规序，(6.1-9) 末项应为 $\chi_i(z)a^*_{\alpha_i}(z)$（对照本节 (6.1-10) 的 $b_i(z)a^*_{\alpha_i}(z)$），下标 $\alpha$ 应为 $\alpha_i$，且指数 $n$ 无处可来；(6.1-8) 两个因子都带星号则 $h_i$ 不再保持权分次。落点：\S6.1.1 式 (6.1-6)--(6.1-9)、原书 p.170。来源：原书 p.170 与 p.168 定理 6.1.3 的对照；并见本译本 (6.1-10)。以上三处均按原书照录，未改公式。}
% [Tutor Note: 式 (6.1-8) 第二个正规序因子印作 $a^*_\beta(z)a^*_\beta(z)$（两个因子都带星号），而定理 6.1.3 中同一表达式印作 $:a^*_\beta(z)a_\beta(z):$。后者才是数算子型的正确形式，此处按原文照录，并在报告中登记为疑似笔误（若照本式理解，$h_i$ 就不再保持权分次）。]

下面去掉正规序并令 $c_i=0$，$i=1,\dots,\ell$。此时相应的 Fourier 系数
作为 $M_{\mathfrak g}$ 上的线性算子不再有良好定义，但它们是空间
$\mathrm{Fun}\,L\cU$ 上定义良好的线性算子。所得的 $L\fg$-模结构容易描述。
事实上，李代数 $L\fg$ 通过向量场作用在 $\mathrm{Fun}\,L\cU$ 上。更一般地，
对任意 $L\mathfrak b$-模 $R$，我们在张量积
$\mathrm{Fun}\,L\cU_{\mathfrak b}\otimes R$ 上得到 $L\fg$ 的自然作用（这
正是由 $L\mathfrak b$-模 $R$ 诱导的拓扑 $L\fg$-模）。若取 $R$ 为使所有
$f_{i,n}$ 作用为 $0$、$h_{i,n}$ 作用为乘以 $\chi_{i,n}$（$i=1,\dots,\ell$，
$n\in\Z$）的一维表示，则相应的 $\mathrm{Fun}\,L\cU_{\mathfrak b}\otimes R
\simeq\mathrm{Fun}\,L\cU$ 上的 $L\fg$-作用由上述公式的 Fourier 系数给出，
只是去掉了正规序并取 $c_i=0$。因此，若去掉正规序并令 $c_i=0$，这些 Fourier
系数确实满足 $L\fg$ 的交换关系。

恢复正规序后，如 \S5.5.2 所解释的，这些交换关系一般会因双缩并而变形。但由
定理 6.1.3 知道，当恢复正规序并令所有 $\chi_{i,n}=0$ 时，存在数 $c_i$ 使
这些 Fourier 系数满足 $\gghat_{\kappa_c}$ 的交换关系（其中 $\mathbf 1=1$）。
项 (6.1-4)--(6.1-6) 不会在交换子中产生新的双缩并。因此对一切非零的
$\chi_{i,n}$，$\gghat_{\kappa_c}$ 的交换关系也成立。证毕。

\renewcommand{\thecorollary}{6.1.5}
\begin{corollary}
    对每个 $\chi(t)\in\mathfrak h^*(\!(t)\!)$，$M_{\mathfrak g}$ 上有临界水平
    的 $\gghat$-模结构，其作用由公式 (6.1-7)--(6.1-9) 给出。
\end{corollary}

称这些模为\textbf{临界水平的 Wakimoto 模}，记作 $W_{\chi(t)}$。

设 $\pi_0$ 是交换代数 $\C[b_{i,n}]_{i=1,\dots,\ell;n<0}$，其导子 $T$ 由公式
\[
    T\cdot b_{i_1,n_1}\cdots b_{i_m,n_m}
    =-\sum_{j=1}^mn_jb_{i_1,n_1}\cdots b_{i_j,n_j-1}\cdots b_{i_m,n_m}
\]
给出。则 $\pi_0$ 自然是交换顶点代数（见 \S2.2.2）。特别地，有
\[
    Y(b_{i,-1},z)=b_i(z)=\sum_{n<0}b_{i,n}z^{-n-1}.
\]
利用与引理 6.1.4 证明中相同的论证，现在得到推论 6.1.2 的加强形式。

\renewcommand{\thetheorem}{6.1.6}
\begin{theorem}
    存在顶点代数同态
    \[
        w_{\kappa_c}:V_{\kappa_c}(\fg)\to M_{\mathfrak g}\otimes\pi_0,
    \]
    使得
    \begin{equation}
        \begin{aligned}
            e_i(z)&\mapsto a_{\alpha_i}(z)+\sum_{\beta\in\Delta_+}
            :\!P^i_\beta(a^*_\alpha(z))a_\beta(z)\!:,\\
            h_i(z)&\mapsto-\sum_{\beta\in\Delta_+}
            \beta(h_i):\!a^*_\beta(z)a_\beta(z)\!:+b_i(z),\\
            f_i(z)&\mapsto\sum_{\beta\in\Delta_+}
            :\!Q^i_\beta(a^*_\alpha(z))a_\beta(z)\!:
            +c_i\partial_za^*_{\alpha_i}(z)+b_i(z)a^*_{\alpha_i}(z),
        \end{aligned}
        \tag{6.1-10}
    \end{equation}
    其中多项式 $P^i_\beta$、$Q^i_\beta$ 由公式 (5.2-12)--(5.2-14) 引入。
\end{theorem}

于是顶点代数 $M_{\mathfrak g}\otimes\pi_0$ 上的任何模都成为
$V_{\kappa_c}(\fg)$-模，从而是临界水平的 $\gghat$-模。我们不要求该模一定
$\Z$-分次。特别地，对任意 $\chi(t)\in\mathfrak h^*(\!(t)\!)$，有一维
$\pi_0$-模 $\C_{\chi(t)}$，其中 $b_{i,n}$ 作用为乘以 $\chi_{i,n}$。相应的
$\gghat_{\kappa_c}$-模就是上面引入的 Wakimoto 模 $W_{\chi(t)}$。

\section{形变到其他水平}

\setcounter{theorem}{0}

\begin{keypoints}
    \item \textbf{核心思想：} 把临界水平的构造连续形变到任意水平 $\kappa$：同态
    $V_\kappa(\fg)\to M_{\mathfrak g}\otimes\pi_0^{\kappa-\kappa_c}$ 的显式
    公式只比临界情形多一个 $(\kappa-\kappa_c)(e_i,f_i)$ 修正项。
    \item \textbf{数学工具：} Heisenberg 顶点代数 $\pi_0^\kappa$、Segal--Sugawara
    向量的像、共形向量、$\operatorname{Der}\mathcal O$ 与拟共形结构、
    $\operatorname{Aut}\mathcal O$-torsor 的扭变。
    \item \textbf{关键结论：} 定理 6.2.1 与命题 6.2.2（共形结构）、推论 6.2.3
    （拟共形结构）、命题 6.2.4 与 6.2.5（无坐标版本）。
    \item \textbf{记号提醒：}$\kappa_c$ 对应 $k=-2$（$\fsl_2$ 情形）；$\rho_{-2}$
    是 $\rho$ 的 $-2$ 模式；$\pi^\kappa_\lambda$ 是 Fock 表示。
\end{keypoints}

本节把 Wakimoto 模的构造推广到非临界水平。我们证明顶点代数同态
\[
    w_\kappa:V_\kappa(\fg)\to M_{\mathfrak g}\otimes\pi_0^{\kappa-\kappa_c}
\]
的存在性，其中 $\pi_0^{\kappa-\kappa_c}$ 是与 Heisenberg 李代数
$\widehat{\mathfrak h}_{\kappa-\kappa_c}$ 相伴的顶点代数，并详细分析这个同态与
坐标变换群 $\operatorname{Aut}\mathcal O$ 在圆盘上作用的相容性。这将使我们
能给出自由场实现的无坐标解释。

\subsection{顶点代数同态}

与前面一样，用 $\mathfrak h$ 表示 $\fg$ 的 Cartan 子代数。设
$\widehat{\mathfrak h}_\kappa$ 是圈代数 $L\mathfrak h=\mathfrak h(\!(t)\!)$ 的
一维中心扩张，其二上循环取为 $L\fg$ 上对应于内积 $\kappa$ 的二上循环的
限制。则由公式 (1.3-4)，$\widehat{\mathfrak h}_\kappa$ 是 Heisenberg 李代数。
我们考虑这个李代数的一个副本，生成元为 $b_{i,n}$（$i=1,\dots,\ell$，
$n\in\Z$）与 $\mathbf 1$，交换关系为
\begin{equation}
    [b_{i,n},b_{j,m}]=n\kappa(h_i,h_j)\delta_{n,-m}\mathbf 1.
    \tag{6.2-1}
\end{equation}
于是 $b_{i,n}$ 满足与 $h_{i,n}$ 相同的交换关系。\transnote{卡点：(6.2-1) 的中心项是 $n\kappa(h_i,h_j)$，而 (6.2-3) 中 $\fsl_2$ 的同一括号却写成 $2kn$，两者如何统一？背景：水平按 $\kappa=k\kappa_0$ 计，$\kappa_0$ 的归一化使 $\kappa_0(h,h)=2$（见 \S1.3，并与 \S6.3.2 的 $\kappa_c=-\frac12\kappa_{\mathfrak g}$ 配套）；$\fsl_2$ 只有一支 $h$，故 $\kappa(h,h)=2k$，(6.2-1) 即化为 $[b_n,b_m]=2kn\delta_{n,-m}\mathbf 1$。同一水平约定也出现在 (6.2-2) 的形变项 $(\kappa-\kappa_c)(e_i,f_i)$ 中。落点：\S6.2.1 式 (6.2-1)、\S6.2.2 式 (6.2-3)。来源：原书 \S6.2.1--\S6.2.2；并见 \S6.3.2 的 $\kappa_c$ 定义。}设 $\pi_0^\kappa$ 是由
% [Tutor Note: 注意 (6.2-1) 的中心项系数是 $n\kappa(h_i,h_j)$，与 (6.2-3) 的 $\fsl_2$ 公式 $[b_n,b_m]=2kn\delta_{n,-m}\mathbf 1$ 对照时要用 $\kappa=k\kappa_0$、$\kappa_0(h,h)=2$。水平约定在 (6.2-2) 的形变项 $(\kappa-\kappa_c)(e_i,f_i)$ 中再次出现，务必统一。]
$\widehat{\mathfrak h}_\kappa$ 的下列一维表示诱导的 $\widehat{\mathfrak h}_\kappa$-模：
该表示定义在由 $b_{i,n}$（$i=1,\dots,\ell$，$n\ge0$）与 $\mathbf 1$ 张成的
阿贝尔李子代数上，其中 $\mathbf 1$ 作用为恒等，其余生成元作用为 $0$。用
$|0\rangle$ 表示这个模的生成向量。它满足 $b_{i,n}|0\rangle=0$，$n\ge0$。
则 $\pi_0^\kappa$ 具有下述 $\Z_+$-分次顶点代数结构（见 [Frenkel and Ben-Zvi
2004] 的定理 2.3.7）：
\begin{itemize}[leftmargin=2em]
    \item $\Z_+$-分次：
    $\deg b_{i_1,n_1}\cdots b_{i_m,n_m}|0\rangle=-\sum_{i=1}^mn_i$；
    \item 真空向量：$|0\rangle$；
    \item 平移算子：$T|0\rangle=0$，$[T,b_{i,n}]=-nb_{i,n-1}$；
    \item 顶点算子：
    \[
        Y(b_{i,-1}|0\rangle,z)=b_i(z)=\sum_{n\in\Z}b_{i,n}z^{-n-1},
    \]
    \[
        Y(b_{i_1,n_1}\cdots b_{i_m,n_m}|0\rangle,z)
        =\prod_{j=1}^m\frac{1}{(-n_j-1)!}
        :\!\partial_z^{-n_1-1}b_{i_1}(z)\cdots
        \partial_z^{-n_m-1}b_{i_m}(z)\!:.
    \]
\end{itemize}
张量积 $M_{\mathfrak g}\otimes\pi_0^{\kappa-\kappa_c}$ 也获得顶点代数结构。

下述定理把定理 6.1.6 的结果推广到临界水平之外。

\renewcommand{\thetheorem}{6.2.1}
\begin{theorem}
    存在顶点代数同态 $w_\kappa:V_\kappa(\fg)\to
    M_{\mathfrak g}\otimes\pi_0^{\kappa-\kappa_c}$，使得
    \begin{equation}
        \begin{aligned}
            e_i(z)&\mapsto a_{\alpha_i}(z)+\sum_{\beta\in\Delta_+}
            :\!P^i_\beta(a^*_\alpha(z))a_\beta(z)\!:,\\
            h_i(z)&\mapsto-\sum_{\beta\in\Delta_+}
            \beta(h_i):\!a^*_\beta(z)a_\beta(z)\!:+b_i(z),\\
            f_i(z)&\mapsto\sum_{\beta\in\Delta_+}
            :\!Q^i_\beta(a^*_\alpha(z))a_\beta(z)\!:\\
            &\qquad+\left(c_i+(\kappa-\kappa_c)(e_i,f_i)\right)
            \partial_za^*_{\alpha_i}(z)+b_i(z)a^*_{\alpha_i}(z),
        \end{aligned}
        \tag{6.2-2}
    \end{equation}
    其中多项式 $P^i_\beta$、$Q^i_\beta$ 由公式 (5.2-12)--(5.2-14) 引入。
\end{theorem}
% [Tutor Note: 式 (6.2-2) 中唯一的“形变项”是 $(\kappa-\kappa_c)(e_i,f_i)\,\partial_za^*_{\alpha_i}(z)$。在临界水平 $\kappa=\kappa_c$ 它消失，回到定理 6.1.3；这一项正是半无限诱导使水平从 $\kappa-\kappa_c$ 平移到 $\kappa$ 的来源，务必与 (6.1-10) 对照核对。]

\begin{proof}
    记 $\widetilde{\cA}^{\mathfrak g}_{\mathrm{loc}}
    =U(M_{\mathfrak g}\otimes\pi_0^{\kappa-\kappa_c})$。由引理 6.1.1，为证明
    定理，只需证明公式 (6.2-2) 定义了李代数同态
    $\gghat_\kappa\to\widetilde{\cA}^{\mathfrak g}_{\mathrm{loc}}$，且把中心元
    $\mathbf 1$ 送到恒等。

    公式 (6.2-2) 当然定义了线性映射
    $\overline w_\kappa:L\fg\to\widetilde{\cA}^{\mathfrak g}_{\mathrm{loc}}$。
    用公式
    \[
        \omega_\kappa(f,g)=\left[\overline w_\kappa(f),
        \overline w_\kappa(g)\right]-\overline w_\kappa([f,g])
    \]
    定义线性映射 $\omega_\kappa:\bigwedge^2L\fg\to
    \widetilde{\cA}^{\mathfrak g}_{\mathrm{loc}}$。与引理 5.6.2 的证明一样
    显式计算，可知 $\omega_\kappa$ 取值于
    $\cA^{\mathfrak g}_{0,\mathrm{loc}}\subset
    \widetilde{\cA}^{\mathfrak g}_{\mathrm{loc}}$。此外，按 $\overline w_\kappa$
    的构造，对任意 $X\in\cA^{\mathfrak g}_{0,\mathrm{loc}}$ 与 $f\in L\fg$，
    有 $[\overline w_\kappa(f),X]=f\cdot X$，其中右端考虑 $f$ 在
    $L\fg$-模 $\cA^{\mathfrak g}_{0,\mathrm{loc}}$ 上的作用。这立刻说明
    $\omega_\kappa$ 是以 $\cA^{\mathfrak g}_{0,\mathrm{loc}}$ 为系数的
    $L\fg$ 二上循环。按构造 $\omega_\kappa$ 是局部的，即属于
    $C^2_{\mathrm{loc}}(L\fg,\cA^{\mathfrak g}_{0,\mathrm{loc}})$。

    下面计算 $\omega_\kappa$ 在 $\bigwedge^2L\mathfrak h$ 上的限制。引理
    5.6.8 证明中的计算表明
    \[
        \omega_\kappa(h_n,h'_m)
        =n\left(\kappa_c(h,h')+(\kappa-\kappa_c)(h,h')\right)
        =n\kappa(h,h').
    \]
    因此这个限制等于 $L\fg$ 上对应于 $\kappa$ 的 Kac--Moody 二上循环
    $\sigma_\kappa$ 的限制。现在引理 5.6.7 说明二上循环 $\omega_\kappa$ 与
    $i_*(\sigma_\kappa)$ 在上同调意义下等价。我们断言它实际上等于
    $i_*(\sigma_\kappa)$。

    事实上，这两个上循环之差是某个元素
    $\gamma\in C^1_{\mathrm{loc}}(L\fg,\cA^{\mathfrak g}_{0,\mathrm{loc}})$
    的上边缘。定理 6.1.3 之前的讨论说明 $\gamma(e_i(z))=\gamma(h_i(z))=0$，
    且 $\gamma(f_i(z))=c'_i\partial_za^*_{\alpha_i}(z)$，其中 $c'_i\in\C$ 是
    某些常数。为求 $c'_i$，我们计算相应的二上循环 $\omega_\kappa+d\gamma$ 在
    $e_{i,n}$ 与 $f_{i,-n}$ 上的值。求得它等于
    $n\sigma_\kappa(e_{i,n},f_{i,-n})+c'_in$。另一方面，$\gghat$ 中的交换关系
    要求它等于 $n\sigma_\kappa(e_{i,n},f_{i,-n})$。因此对一切
    $i=1,\dots,\ell$ 有 $c'_i=0$，从而 $\gamma=0$。于是
    $\omega_\kappa=i_*(\sigma_\kappa)$。这说明公式 (6.2-2) 确实定义了把中心元
    $\mathbf 1$ 送到恒等的李代数同态
    $\gghat_\kappa\to\cA^{\mathfrak g}_{\le1,\mathrm{loc}}$。证毕。
\end{proof}

现在顶点代数 $M_{\mathfrak g}\otimes\pi_0^{\kappa-\kappa_c}$ 上的任何模都
成为 $V_\kappa(\fg)$-模，从而是 $\gghat_\kappa$-模（$K$ 作用为 1）。对
$\lambda\in\mathfrak h^*$，设 $\pi^{\kappa-\kappa_c}_\lambda$ 是由向量
$|\lambda\rangle$ 生成的 $\widehat{\mathfrak h}_{\kappa-\kappa_c}$ 的 Fock
表示，满足
\[
    b_{i,n}|\lambda\rangle=0,\quad n>0;\qquad
    b_{i,0}|\lambda\rangle=\lambda(h_i)|\lambda\rangle;\qquad
    \mathbf 1|\lambda\rangle=|\lambda\rangle.
\]
则
\[
    W_{\lambda,\kappa}\overset{\mathrm{def}}{=}
    M_{\mathfrak g}\otimes\pi^{\kappa-\kappa_c}_\lambda
\]
是 $M_{\mathfrak g}\otimes\pi_0^{\kappa-\kappa_c}$-模，从而是 $\gghat$-模。
称之为\textbf{水平 $\kappa$、最高权 $\lambda$ 的 Wakimoto 模}。

\subsection{$\widehat{\fsl}_2$ 上的 Wakimoto 模}

本节显式描述 $\widehat{\fsl}_2$ 上的 Wakimoto 模。

设 $\{e,h,f\}$ 是李代数 $\fsl_2$ 的标准基。设 $\kappa_0$ 是 $\fsl_2$ 上按
$\kappa_0(h,h)=2$ 归一化的不变内积。把 $\fsl_2$ 上任一不变内积 $\kappa$ 写成
$k\kappa_0$，其中 $k\in\C$，并在记号中用 $k$ 代替 $\kappa$。特别地，临界
内积 $\kappa_c$ 对应 $k=-2$。$\fsl_2$ 情形的 $\Delta_+$ 只有一个元素，因此
我们省略 $a_\alpha(z)$ 与 $a^*_\alpha(z)$ 中的指标 $\alpha$。同样省略
$b_i(z)$ 等中的指标 $i$，并把 $M_{\mathfrak{sl}_2}$ 记为 $M$。我们还通过
$\chi\mapsto\chi(h)$ 把 Cartan 子代数的对偶空间 $\mathfrak h^*$ 与 $\C$
等同。

Weyl 代数 $\cA_{\mathfrak{sl}_2}$ 的生成元为 $a_n,a^*_n$（$n\in\Z$），交换
关系为
\[
    [a_n,a^*_m]=\delta_{n,-m}.
\]
其 Fock 表示记为 $M$。Heisenberg 李代数 $\widehat{\mathfrak h}_k$ 的生成元为
$b_n$（$n\in\Z$）与 $\mathbf 1$，交换关系为
\[
    [b_n,b_m]=2kn\delta_{n,-m}\mathbf 1,
\]
$\pi^k_\lambda$ 是由向量 $|\lambda\rangle$ 生成的 Fock 表示，满足
\[
    b_n|\lambda\rangle=0,\quad n>0;\qquad
    b_0|\lambda\rangle=\lambda|\lambda\rangle;\qquad
    \mathbf 1|\lambda\rangle=|\lambda\rangle.
\]
模 $\pi_0^k$ 与张量积 $M\otimes\pi_0^k$ 都是顶点代数。

顶点代数同态 $w_k:V_k(\fsl_2)\to M\otimes\pi_0^{k+2}$ 由下列公式给出：
\begin{equation}
    \begin{aligned}
        e(z)&\mapsto a(z),\\
        h(z)&\mapsto-2:\!a^*(z)a(z)\!:+b(z),\\
        f(z)&\mapsto-:\!a^*(z)^2a(z)\!:+k\partial_za^*(z)+a^*(z)b(z).
    \end{aligned}
    \tag{6.2-3}
\end{equation}
Wakimoto 模 $M\otimes\pi^{k+2}_\lambda$ 将记作 $W_{\lambda,k}$，其最高权向量
记作 $|\lambda\rangle$。\transnote{卡点：为何 $\fsl_2$ 的临界水平恰好是 $k=-2$？背景：临界水平由 $\kappa_c=-\frac12\kappa_{\mathfrak g}$ 定义（见 \S6.3.2）；对 $\fsl_2$ 有 $\kappa_{\mathfrak g}=4\kappa_0$ 且 $\kappa_0(h,h)=2$，故 $\kappa_c=-2\kappa_0$，即 $k=-2$。此时 $W_{\lambda,k}=M\otimes\pi^{k+2}_\lambda$ 中 Heisenberg 的水平 $k+2=0$，$\pi_0$ 退化为交换顶点代数，这正是 \S6.2.4 讨论临界拟共形结构的出发点。落点：\S6.2.2、\S6.3.2。来源：原书 \S6.3.2 的 $\kappa_c$ 定义与 \S6.2.2 的式 (6.2-3)；$\kappa_{\mathfrak g}=4\kappa_0$ 由 $\fsl_2$ 的 Killing 形式直接算出。}
% [Tutor Note: $k=-2$ 是 $\fsl_2$ 的临界水平，此时 (6.2-3) 中的 $k\partial_za^*$ 项变为 $-2\partial_za^*$；与一般公式 (6.2-2) 中的修正项 $(\kappa-\kappa_c)(e,f)\partial_za^*$ 对照时，注意 $\kappa_c(e,f)$ 的符号与归一化（本处 $\kappa_c$ 对应 $k=-2$）。这是核对符号约定的关键处。]

我们将在第 7 章使用这些显式公式构造 Wakimoto 模之间的交织算子。

\subsection{非临界水平的共形结构}

本节证明：当 $\kappa\ne\kappa_c$ 时，定理 6.2.1 的同态 $w_\kappa$ 是共形顶点
代数同态（共形顶点代数的定义见 \S3.1.4）。这将使我们得到该同态的无坐标
版本。取极限 $\kappa\to\kappa_c$，我们还将得到 $w_{\kappa_c}$ 的无坐标版本。

顶点代数 $V_\kappa(\fg)$（$\kappa\ne\kappa_c$）具有由 Segal--Sugawara 向量
\begin{equation}
    \mathbf S_\kappa=\frac12\sum_aJ^a_{-1}J_{a,-1}|0\rangle
    \tag{6.2-4}
\end{equation}
给出的共形顶点代数结构，其中 $\{J_a\}$ 是 $\fg$ 的基，关于内积
$\kappa-\kappa_c$ 与基 $\{J^a\}$ 对偶（见 \S3.1.4，其中这个向量记作
$\widetilde S$）。我们需要计算 $\mathbf S_\kappa$ 在 $w_\kappa$ 下的像。

\renewcommand{\theproposition}{6.2.2}
\begin{proposition}
    $\mathbf S_\kappa$ 在 $w_\kappa$ 下的像等于
    \begin{equation}
        \left(\sum_{\alpha\in\Delta_+}a_{\alpha,-1}a^*_{\alpha,-1}
        +\frac12\sum_{i=1}^\ell b_{i,-1}b^i_{-1}-\rho_{-2}\right)|0\rangle,
        \tag{6.2-5}
    \end{equation}
    其中 $\{b^i\}$ 是 $\{b_i\}$ 的对偶基，$\rho$ 是在内积
    $(\kappa-\kappa_c)|_{\mathfrak h}$ 诱导的同构下对应于 $\rho\in\mathfrak h^*$
    的 $\mathfrak h$ 中元素。\transnote{卡点：(6.2-5) 末项的 $-\rho_{-2}$ 是什么，为什么不能省？背景：此处的 $\rho$ 指由内积 $(\kappa-\kappa_c)|_{\mathfrak h}$ 把 $\rho\in\mathfrak h^*$ 对应的 $\mathfrak h$ 中元素，$\rho_{-2}$ 是相应 Heisenberg 场的 $-2$ 模式，故 $\rho_{-2}|0\rangle$ 与 (6.2-6) 中的 $b_{i,-2}|0\rangle$ 同处次数 $2$、权 $0$ 的齐次分量。它是补上 $h_i$ 项里 $a^*a$ 双缩并、使 $w_\kappa(\mathbf S_\kappa)$ 真正满足共形向量公理的必需修正；同一项的残余在 (6.2-13) 第三式给出常数修正（与 $\langle\rho,h_i\rangle=1$ 有关）。落点：\S6.2.3 命题 6.2.2、\S6.2.4 式 (6.2-13)、原书 p.175。来源：原书 p.175 式 (6.2-5) 及其证明；并见 (6.2-12) 与 \S6.2.4。}
% [Tutor Note: (6.2-5) 中 $-\rho_{-2}$ 项不是“零模”：它是 $\rho\in\mathfrak h$ 的 $-2$ 模式，来自 (6.2-11) 中 $h_i$ 项的双缩并修正（与 $\langle\rho,h_i\rangle=1$ 有关）。缺少 $-\rho_{-2}$ 时 $w_\kappa(\mathbf S_\kappa)$ 不再是共形向量——这是核对共形结构时最容易漏掉的一项。]
\end{proposition}

向量 $w_\kappa(\mathbf S_\kappa)$ 定义了共形代数结构，从而是
$M_{\mathfrak g}\otimes\pi_0^{\kappa-\kappa_0}$ 上的 Virasoro 代数作用。
同态 $w_\kappa$ 把 $M_{\mathfrak g}\otimes\pi_0^{\kappa-\kappa_0}$ 上相应的
$(\operatorname{Der}\mathcal O,\operatorname{Aut}\mathcal O)$ 作用与
$V_\kappa(\fg)$ 上的自然
$(\operatorname{Der}\mathcal O,\operatorname{Aut}\mathcal O)$ 作用交织起来。\transnote{卡点：本段两处印作 $\pi_0^{\kappa-\kappa_0}$，全章其余各处均作 $\pi_0^{\kappa-\kappa_c}$，应取哪一个？背景：$\pi_0^\nu$ 的上标是 Heisenberg 代数 $\widehat{\mathfrak h}_\nu$ 的水平，而 Wakimoto 构造把 $\fg$ 的水平 $\kappa$ 与 $\mathfrak h$ 的水平 $\kappa-\kappa_c$ 配对，故此处只能是 $\kappa_c$；紧随其后的证明起首即写作 $M_{\mathfrak g}\otimes\pi_0^{\kappa-\kappa_c}$，同页自证。落点：\S6.2.3 命题 6.2.2 之后一段、原书 p.175。来源：原书 p.175（题设段与其证明的对照）；并见 \S6.2.1 与式 (6.2-2)。此为原书排印之误，译文照录，未改公式。}
% [Tutor Note: 此处原书印作 $\pi_0^{\kappa-\kappa_0}$（下标 $0$），而全章其余处均为 $\pi_0^{\kappa-\kappa_c}$；按上下文应为 $\kappa_c$。按原文照录，并在报告中登记为疑似笔误。]

\begin{proof}
    顶点代数 $V_\kappa(\fg)$ 与 $M_{\mathfrak g}\otimes\pi_0^{\kappa-\kappa_c}$
    都带有 $\Z$-分次，以及由 $\mathfrak h$ 作用而来的、关于 $\fg$ 的根格的
    权重分次（注意 $\mathrm{wt}\,a_{\alpha,n}=-\mathrm{wt}\,a^*_{\alpha,n}
    =\alpha$，$\mathrm{wt}\,b_{i,n}=0$）。同态 $w_\kappa$ 保持这些分次。向量
    $\mathbf S_\kappa\in V_\kappa(\fg)$ 关于 $\Z$-分次的次数为 $2$，关于根格
    分次的权为 $0$。因此 $w_\kappa(\mathbf S_\kappa)$ 也有同样的性质。
    $M_{\mathfrak g}\otimes\pi_0^{\kappa-\kappa_c}$ 相应齐次子空间的基由形如
    \begin{equation}
        b_{i,-1}b_{j,-1},\qquad b_{i,-2},\qquad
        a_{\alpha,-1}a^*_{\alpha,-1},
        \tag{6.2-6}
    \end{equation}
    \begin{equation}
        a_{\alpha,-1}a_{\beta,-1}a^*_{\alpha,0}a^*_{\beta,0},\qquad
        a_{\alpha,-1}a^*_{\alpha,0}b_{i,-1},
        \tag{6.2-7}
    \end{equation}
    \begin{equation}
        a_{\alpha,-2}a^*_{\alpha,0},\qquad
        a_{\alpha+\beta,-2}a^*_{\alpha,0}a^*_{\beta,0},\qquad
        a_{\alpha+\beta,-1}a^*_{\alpha,-1}a^*_{\beta,0},
        \tag{6.2-8}
    \end{equation}
    的单项式作用到真空向量 $|0\rangle$ 上得到。我们要证明
    $w_\kappa(\mathbf S_\kappa)$ 是 (6.2-6) 中单项式的线性组合。

    按构造，顶点算子 $w_\kappa(\mathbf S_\kappa)$ 的 Fourier 系数
    $L_n$（$n\in\Z$）保持 $M_{\mathfrak g}\otimes\pi_0^{\kappa-\kappa_c}$
    上的权重分次，并且关于顶点代数分次有 $\deg L_n=-n$。我们断言：形如
    $P(a^*_{\alpha,0})|0\rangle$ 的任意向量都被 $L_n$（$n\ge0$）湮灭。
    对 $n>0$ 这由次数理由显然。为看出 $L_0$ 也如此，注意
    \[
        L_0\cdot P(a^*_{\alpha,0})|0\rangle
        =\frac12\sum_aw_\kappa(J_{a,0})w_\kappa(J^a_0)
        \cdot P(a^*_{\alpha,0})|0\rangle.
    \]
    按定理 6.2.1 中同态 $w_\kappa$ 的公式，常数子代数
    $\fg\subset\gghat_\kappa$ 通过 $w_\kappa$ 在
    $\C[a^*_{\alpha,0}]_{\alpha\in\Delta_+}|0\rangle$ 上的作用，与 $\fg$ 在
    $\C[y_\alpha]_{\alpha\in\Delta_+}=\mathrm{Fun}\,N_+$ 上的自然作用一致，
    只要作代换 $a^*_{\alpha,0}\mapsto y_\alpha$。因此 $L_0$ 在
    $P(a^*_{\alpha,0})|0\rangle$ 上的作用与 Casimir 算子
    $\frac12\sum_aJ_aJ^a$ 在此代换下在
    $\C[y_\alpha]_{\alpha\in\Delta_+}$ 上的作用一致。但作为 $\fg$-模，后者
    正是我们在 \S5.2.3 中证明的反变 Verma 模 $M^*_0$。因此 Casimir 算子在
    其上的作用等于 $0$。

    形如 $P(a^*_{\alpha,0})|0\rangle$ 的所有向量都被 $L_n$（$n\ge0$）湮灭，
    这一事实排除了单项式 (6.2-7) 与 (6.2-8) 出现在
    $w_\kappa(\mathbf S_\kappa)$ 中（(6.2-8) 的最后一个单项式除外）。为消去
    (6.2-8) 中最后一个单项式，我们证明
    \[
        L_n\cdot a_{\alpha,-1}|0\rangle=0,\quad n>0;
        \qquad
        L_0\cdot a_{\alpha,-1}|0\rangle=a_{\alpha,-1}|0\rangle,
        \quad\alpha\in\Delta_+.
    \]
    第一式由次数理由成立。

    为证第二式，设 $\{e_\alpha\}_{\alpha\in\Delta_+}$ 是
    $\mathfrak n_+\subset\fg$ 的根基，使 $e_{\alpha_i}=e_i$。则向量
    $e_{\alpha,-1}|0\rangle\in V_\kappa(\fg)$，从而
    $w_\kappa(e_{\alpha,-1}|0\rangle)\in
    M_{\mathfrak g}\otimes\pi_0^{\kappa-\kappa_c}$ 被 $L_n$（$n>0$）湮灭，
    并且是 $L_0$ 的特征向量，特征值为 $1$。我们有
    \begin{equation}
        w_\kappa(e_{\alpha,-1}|0\rangle)=a_{\alpha,-1}|0\rangle
        +\sum_{\beta\in\Delta_+;\beta>\alpha}
        P^\alpha_\beta(a^*_{\alpha,0})a_{\beta,-1}|0\rangle,
        \tag{6.2-9}
    \end{equation}
    其中多项式 $P^\alpha_\beta$ 由 $e_\alpha$ 在 $\cU$ 上作用的下列公式确定：
    \[
        e_\alpha=\frac{\partial}{\partial y_\alpha}
        +\sum_{\beta\in\Delta_+;\beta>\alpha}
        P^\alpha_\beta(y_\alpha)\frac{\partial}{\partial y_\beta}
    \]
    （见公式 (6.1-2)）。

    从最大根 $\alpha_{\max}$ 开始，按根的高递减作归纳，我们由这些公式推出
    $L_0\cdot a_{\alpha,-1}|0\rangle=a_{\alpha,-1}|0\rangle$。这消去了
    (6.2-8) 的最后一个单项式，并给出公式
    \begin{equation}
        w_\kappa(\mathbf S_\kappa)
        =\sum_{\alpha\in\Delta_+}a_{\alpha,-1}a^*_{\alpha,-1}|0\rangle
        \tag{6.2-10}
    \end{equation}
    再加上 (6.2-6) 中前两类单项式之和。剩下的问题是确定它们进入公式的系数。

    为此使用 $w_\kappa(h_{i,-1}|0\rangle)$ 的下列公式，它由定理 6.2.1 推出：
    \begin{equation}
        w_\kappa(h_{i,-1}|0\rangle)
        =\left(-\sum_{\beta\in\Delta_+}
        \beta(h_i)a^*_{0,\beta}a_{\beta,-1}+b_{i,-1}\right)|0\rangle.
        \tag{6.2-11}
    \end{equation}
    由它求出 $L_0$ 与 $L_1$ 在 $w_\kappa(h_{i,-1}|0\rangle)$ 第一个求和项上的
    作用：
    \[
        L_0\cdot\left(-\sum_{\beta\in\Delta_+}
        \beta(h_i)a^*_{0,\beta}a_{\beta,-1}|0\rangle\right)
        =-\sum_{\beta\in\Delta_+}
        \beta(h_i)a^*_{0,\beta}a_{\beta,-1}|0\rangle,
    \]
    \[
        L_1\cdot\left(-\sum_{\beta\in\Delta_+}
        \beta(h_i)a^*_{0,\beta}a_{\beta,-1}|0\rangle\right)
        =-\sum_{\beta\in\Delta_+}\beta(h_i)|0\rangle
        =-2\rho(h_i)|0\rangle.
    \]
    此外，$L_n$（$n>1$）在它上面作用为 $0$。

    另一方面，我们知道 $w_\kappa(h_{i,-1}|0\rangle)$ 必须被 $L_n$（$n>0$）
    湮灭，并且是 $L_0$ 的特征向量，特征值为 $1$。与 (6.2-10) 相加后满足这个
    条件的、(6.2-6) 前两类单项式的唯一组合是
    \begin{equation}
        \frac12\sum_{i=1}^\ell b_{i,-1}b^i_{-1}-\rho_{-2}.
        \tag{6.2-12}
    \end{equation}
    这就证明了公式 (6.2-5)。现在利用引理 3.1.2 直接验证：由该公式给出的向量
    $w_\kappa(\mathbf S_\kappa)$ 满足共形向量的公理。证毕。
\end{proof}

\subsection{临界水平的拟共形结构}

现在利用上述引理获得关于同态 $w_{\kappa_c}$ 的更多信息。用 $\mathcal O$
表示完备拓扑环 $\C[[t]]$，用 $\operatorname{Der}\mathcal O$ 表示其连续导子
构成的李代数。注意
$\operatorname{Der}\mathcal O\simeq\C[[t]]\partial_t$。

回忆：若顶点代数 $V$ 带有李代数 $\operatorname{Der}\mathcal O$ 的作用，满足
下列条件，则称 $V$ 是\textbf{拟共形的}：
\begin{itemize}[leftmargin=2em]
    \item 对一切 $A\in V$ 有公式
    \[
        [L_m,A_{(k)}]=\sum_{n\ge-1}\binom{m+1}{n+1}
        (L_n\cdot A)_{(m+k-n)};
    \]
    \item 元素 $L_{-1}=-\partial_t$ 作为平移算子 $T$ 作用；
    \item 元素 $L_0=-t\partial_t$ 半单地作用，特征值为整数；
    \item 李子代数 $\operatorname{Der}_+\mathcal O$ 局部幂零地作用
\end{itemize}
（见 [Frenkel and Ben-Zvi 2004] 的定义 6.3.4）。特别地，共形顶点代数自动是
拟共形的（$\operatorname{Der}\mathcal O$ 的作用来自 Virasoro 作用）。

李代数 $\operatorname{Der}\mathcal O$ 自然地作用在 $\gghat_{\kappa_c}$ 上
并保持 $\fg[[t]]$，从而作用在 $V_{\kappa_c}(\fg)$ 上。
$V_{\kappa_c}(\fg)$ 上的 $\operatorname{Der}\mathcal O$ 作用等于
$V_\kappa(\fg)$（$\kappa\ne\kappa_c$）上由 Sugawara 共形结构得到的
$\operatorname{Der}\mathcal O$ 作用当 $\kappa\to\kappa_c$ 时的极限。因此
这个作用在 $V_{\kappa_c}(\fg)$ 上定义了拟共形顶点代数结构。

下面如下定义 $M_{\mathfrak g}\otimes\pi_0$ 上的拟共形代数结构。顶点代数
$M_{\mathfrak g}$ 以 (6.2-10) 为共形向量，因此是共形的，从而是拟共形的。
交换顶点代数 $\pi_0$ 是以 (6.2-12) 为共形向量的共形顶点代数族
$\pi_0^{\kappa-\kappa_c}$ 当 $\kappa\to\kappa_c$ 时的极限。李代数
$\operatorname{Der}\mathcal O$ 在 $\pi_0^{\kappa-\kappa_c}$ 上的诱导作用在
极限 $\kappa\to\kappa_c$ 时有良好定义，因此诱导出 $\pi_0$ 上的
$\operatorname{Der}\mathcal O$ 作用。于是它给出 $\pi_0$ 上的拟共形顶点代数
结构。事实上，$\operatorname{Der}\mathcal O$ 作用由代数结构
$\pi_0\simeq\C[b_{i,n}]$ 的导子给出，因此由 [Frenkel and Ben-Zvi 2004] 的
引理 6.3.5，它在 $\pi_0$ 上定义了拟共形顶点代数结构。具体地说，
$\operatorname{Der}\mathcal O$ 的基元素 $L_n=-t^{n+1}\partial_t$（$n\ge-1$）
在 $\pi_0$ 上的作用由下列公式确定：
\begin{equation}
    \begin{aligned}
        L_n\cdot b_{i,m}&=-mb_{i,n+m}, & &-1\le n<-m,\\
        L_n\cdot b_{i,-n}&=n(n+1), & &n>0,\\
        L_n\cdot b_{i,m}&=0, & &n>-m
    \end{aligned}
% [Tutor Note: (6.2-13) 第三式 $L_n\cdot b_{i,-n}=n(n+1)$ 是常系数修正，来自共形向量中的 $-\rho_{-2}$（见 (6.2-12)）在极限下的残余，正比于 $\langle\rho,h_i\rangle=1$。若改用 $\langle\rho,h_i\rangle\ne1$ 的归一化，这一式的常数会随之改变。]
    \tag{6.2-13}
\end{equation}
（注意对一切 $i$ 有 $\langle\rho,h_i\rangle=1$）。现在把上述
$\operatorname{Der}\mathcal O$ 作用相加，得到 $M_{\mathfrak g}\otimes\pi_0$
上的拟共形顶点代数结构。

由于 $V_{\kappa_c}(\fg)$ 与 $M_{\mathfrak g}\otimes\pi_0$ 上的拟共形结构都
来自共形结构当 $\kappa\to\kappa_c$ 的极限，我们得到命题 6.2.2 的如下推论。

\renewcommand{\thecorollary}{6.2.3}
\begin{corollary}
    同态 $w_{\kappa_c}:V_{\kappa_c}(\fg)\to M_{\mathfrak g}\otimes\pi_0$
    保持拟共形结构。特别地，它把两边的 $\operatorname{Der}\mathcal O$ 与
    $\operatorname{Aut}\mathcal O$ 作用交织起来。
\end{corollary}

\subsection{场的变换公式}

现在可以得到场 $a_\alpha(z)$、$a^*_\alpha(z)$ 与 $b_i(z)$ 以及模
$M_{\mathfrak g}$ 与 $\pi_0$ 的变换公式。

由共形向量 $w_\kappa(\mathbf S_\kappa)$ 的显式公式 (6.2-5)，得到下列交换
关系：
\[
    [L_n,a_{\alpha,m}]=-ma_{m+n},\qquad
    [L_n,a^*_{\alpha,m}]=-(m-1)a^*_{\alpha,m-1}.
\]
与 \S3.5.5 一样，由此推出 $a_\alpha(z)$ 在去心圆盘
$D^\times=\Spec\C(\!(z)\!)$ 上按一次形式变换，而 $a^*_\alpha(z)$ 按函数
变换。特别地，得到模 $M_{\mathfrak g}$ 的如下描述。\transnote{卡点：这两条交换关系与紧接着的"$a_\alpha(z)$ 按一次形式、$a^*_\alpha(z)$ 按函数变换"相互矛盾。背景：共形权为 $1$ 的场应满足 $[L_n,a_{\alpha,m}]=-(m+n)a_{\alpha,m+n}$，共形权为 $0$ 的场应满足 $[L_n,a^*_{\alpha,m}]=-(m+n)a^*_{\alpha,m+n}$；原书印作 $-ma_{m+n}$ 与 $-(m-1)a^*_{\alpha,m-1}$，前者漏了下标 $\alpha$ 与 $n$，后者指数也不对。落点：\S6.2.5、原书 p.179；一般公式见 \S3.5.5。来源：原书 p.179 同页下文的"一次形式/函数"陈述与 \S3.5.5 的一般公式。按公式零改动原则，此处保留原式，仅作说明。}
% [Tutor Note: 这两个交换关系照原书印刷：第一式右端印作 $-ma_{m+n}$（未标 $\alpha$），第二式印作 $-(m-1)a^*_{\alpha,m-1}$。按场的共形权应有 $[L_n,a_{\alpha,m}]=-(m+n)a_{\alpha,m+n}$ 与 $[L_n,a^*_{\alpha,m}]=-(m+n)a^*_{\alpha,m+n}$（$a$ 的权为 1、$a^*$ 的权为 0）；原书这两个公式疑有排印错误。译文照录原文，请核对原书 p.179 与 \S3.5.5 的一般公式。]

把 $\cU$ 与 $\mathfrak n_+$ 等同。考虑 Heisenberg 李代数 $\Gamma$，它是
交换李代数 $\mathfrak n_+\otimes\cK\oplus\mathfrak n^*_+\otimes
\Omega_\cK$ 的中心扩张，上循环由公式
\[
    f(t),g(t)dt\mapsto\int\langle f(t),g(t)\rangle dt
\]
给出。这个上循环与坐标无关，因此 $\Gamma$ 带有
$\operatorname{Der}\mathcal O$ 的自然作用，它保持李子代数
$\Gamma_+=\mathfrak n_+\otimes\mathcal O\oplus\mathfrak n^*_+\otimes
\Omega_{\mathcal O}$。把完备 Weyl 代数
$\widetilde{\cA}^{\mathfrak g}$ 与 $U(\Gamma)/(\mathbf 1-1)$ 的某个完备化
等同，其中 $\mathbf 1$ 是中心元。于是模 $M_{\mathfrak g}$ 与由
$\Gamma_+\oplus\C\mathbf 1$ 的一维表示诱导的 $\Gamma$-模等同：其中
$\Gamma_+$ 作用为 $0$，$\mathbf 1$ 作用为恒等。上面考虑的 $M_{\mathfrak g}$
上的 $\operatorname{Der}\mathcal O$ 作用正是这个诱导模上的自然作用。

现在考虑场 $b_i(z)=\sum_{n<0}b_{i,n}z^{-n-1}$ 与模 $\pi_0$。公式 (6.2-13)
描述了 $\operatorname{Der}\mathcal O$ 在 $b_{i,n}$ 上从而在级数 $b_i(z)$ 上
的作用。事实上，这些公式说明 $\partial_z+b_i(z)$ 按线丛
$\Omega^{-\langle\rho,h_i\rangle}$ 上的联络变换。

更确切地说，设 $^LH$ 是 $H$ 的对偶群，即满足下述性质的复环面：其特征标格
$^LH\to\C^\times$ 就是 $H$ 的余特征标格（即同态 $\C^\times\to H$），而
$^LH$ 的余特征标格就是 $H$ 的特征标格。\footnote{因此 $^LH$ 是 \S1.1.5 引入的 $G$ 的 Langlands
对偶群 $^LG$ 的 Cartan 子群。} 于是 $^LH$ 的李代数 $L\mathfrak h$ 自然地
与 $\mathfrak h^*$ 等同。用 $\Omega^{-\rho}$ 表示圆盘 $D$ 上唯一的主
$^LH$-丛，使得与任一特征标 $\check\lambda:{}^LH\to\C^\times$（等价地，$H$
的余特征标）相伴的线丛为
$\Omega^{-\langle\rho,\check\lambda\rangle}$（见 \S4.2.3）。用
$\mathrm{Conn}(\Omega^{-\rho})_D$ 表示这个 $^LH$-丛上所有联络构成的空间。
它是 $\mathfrak h^*\otimes\Omega_{\mathcal O}$ 上的一个 torsor。

关于 $b_i(z)$ 的上述论断可以重述如下：考虑 $\mathfrak h^*$-值场
$\mathbf b(z)=\sum_{i=1}^\ell b_i(z)\omega_i$，其中
$\langle\mathbf b(z),h_i\rangle=b_i(z)$。则算子
$\partial_z+\mathbf b(z)$ 按 $D$ 上 $^LH$-丛 $\Omega^{-\rho}$ 上的联络变换。
等价地，$\partial_z+b_i(z)$ 按 $D$ 上线丛
$\Omega^{-\langle\rho,h_i\rangle}$ 上的联络变换。

为看出这一点，设 $w$ 是新坐标，使得 $z=\varphi(w)$；则同一联络将表现为
$\partial_w+\widetilde{\mathbf b}(w)$，其中
\begin{equation}
    \widetilde{\mathbf b}(w)
    =\varphi'\cdot\mathbf b(\varphi(w))
    +\rho\frac{\varphi''}{\varphi'}.
    \tag{6.2-14}
\end{equation}
容易验证公式 (6.2-14) 与 (6.2-13) 等价。

这说明 $\pi_0$ 同构于函数代数
$\mathrm{Fun}(\mathrm{Conn}(\Omega^{-\rho})_D)$。若在 $D$ 上选取坐标 $z$，
则把 $\mathrm{Conn}(\Omega^{-\rho})_D$ 与
$\mathfrak h^*\otimes\Omega_{\mathcal O}$ 等同，而把这个代数与自由多项式
代数 $\C[b_{i,n}]_{i=1,\dots,\ell;n<0}$ 等同，其中生成元 $b_{i,n}$ 是
$\mathfrak h^*\otimes\Omega_{\mathcal O}\simeq\mathfrak h^*[[z]]dz$ 上的
下列线性泛函：
\[
    b_{i,n}(\chi(z)dz)=\Res_{z=0}\langle\chi(z),h_i\rangle z^n dz.
\]

\subsection{无坐标版本}

迄今为止，我们把 Wakimoto 模看作李代数 $\gghat_{\kappa_c}$ 的表示，后者是
$L\fg=\fg(\!(t)\!)$ 的中心扩张。正如 \S3.5.2 所解释的，重要的是发展一套
适用于李代数 $\fg(\cK_x)=\fg\otimes\cK_x$ 的中心扩张
$\gghat_{\kappa_c,x}$ 的理论，其中 $\cK_x$ 是光滑曲线上点 $x$ 处去心圆盘上
的函数代数。换言之，$\cK_x$ 是 $X$ 上对应于 $x$ 的函数域在 $x$ 处的完备化。
它是拓扑代数，同构于 $\C(\!(t)\!)$，但同构不是典范的。若在 $x$ 处选取形式
坐标 $t$，可以把 $\cK_x$ 与 $\C(\!(t)\!)$ 等同，但这个等同不是典范的，因为
在 $x$ 处通常没有优先的坐标 $t$。然而，如 \S3.5.2 所见，我们有
$\fg(\cK_x)$ 的典范中心扩张 $\gghat_{\kappa_c,x}$ 与真空
$\gghat_{\kappa_c,x}$-模
\[
    V_\kappa(\fg)_x
    =\Ind_{\fg\otimes\mathcal O_x\oplus\C\mathbf 1}^{\gghat_{\kappa_c,x}}\C,
\]
其中如前，$\mathcal O_x$ 表示 $\cK_x$ 中的整数环，它同构于 $\C[[t]]$。

正如 \S3.5.4 所解释的，$V_\kappa(\fg)_x$ 也可以通过 $V_\kappa(\fg)$ 用
$x$ 处形式坐标的 $\operatorname{Aut}\mathcal O$-torsor
$\mathcal Aut_x$ 作扭变得到：
\[
    V_\kappa(\fg)_x=\mathcal Aut_x\times_{\operatorname{Aut}\mathcal O}
    V_\kappa(\fg).
\]
现在希望把定理 6.2.1 的自由场实现同态
$w_\kappa:V_\kappa(\fg)\to M_{\mathfrak g}\otimes\pi_0^{\kappa-\kappa_c}$
用无坐标方式重述。

对于点 $x$，有上一节引入的 Heisenberg 李代数 $\Gamma$ 的类似物。按定义，
这个李代数（记作 $\Gamma_x$）是
$(\mathfrak n_+\otimes\cK_x)\oplus(\mathfrak n^*_+\otimes\Omega_{\cK_x})$
的中心扩张。设 $\Gamma_{+,x}$ 是其交换李子代数
$(\mathfrak n_+\otimes\mathcal O_x)\oplus(\mathfrak n^*_+\otimes
\Omega_{\mathcal O_x})$。设 $M_{\mathfrak g,x}$ 是 $\Gamma_x$-模
\[
    M_{\mathfrak g,x}=\Ind_{\Gamma_{+,x}\oplus\C\mathbf 1}^{\Gamma_x}\C.
\]
与 $V_\kappa(\fg)_x$ 的情形一样，有
\[
    M_{\mathfrak g,x}=\mathcal Aut_x\times_{\operatorname{Aut}\mathcal O}
    M_{\mathfrak g},
\]
其中 $\operatorname{Aut}\mathcal O$ 在 $M_{\mathfrak g}$ 上的作用来自坐标
变换经其在李代数 $\Gamma$ 上的作用所诱导的 $M_{\mathfrak g}$ 上的作用，如
上一节所述。

其次，需要 Heisenberg 李代数 $\mathfrak h_\nu$ 及其模 $\pi^\nu_0$ 的、附着
于点 $x$ 的版本。

考虑 $^LH$-丛 $\Omega^{-\rho}$ 在
$D^\times_x=\Spec\cK_x$ 上的所有可能的复数值 $\lambda$ 所对应的
$\lambda$-联络构成的空间
$\mathrm{Conn}_{\{\lambda\}}(\Omega^{-\rho})_{D^\times_x}$。若选取同构
$\cK_x\simeq\C(\!(t)\!)$，则 $\lambda$-联络是算子
$\nabla=\lambda\partial_t+\chi(t)$，其中 $\chi(t)\in\mathfrak h^*(\!(t)\!)$。
有正合列
\[
    0\to\mathfrak h^*\otimes\Omega_{\cK_x}
    \to\mathrm{Conn}_{\{\lambda\}}(\Omega^{-\rho})_{D^\times_x}
    \to\C\partial_t\to0,
\]
其中倒数第二个映射把如上 $\nabla$ 送到 $\lambda\partial_t$。

设 $\widehat{\mathfrak h}_{\nu,x}$ 是
$\mathrm{Conn}_{\{\lambda\}}(\Omega^{-\rho})_{D^\times_x}$ 的拓扑对偶向量
空间。它适合正合列（这里使用 $\Omega_{\cK_x}$ 与 $\cK_x$ 之间的留数配对）
\[
    0\to\C\mathbf 1\to\widehat{\mathfrak h}_{\nu,x}
    \to\mathfrak h\otimes\cK_x\to0,
\]
其中 $\mathbf 1$ 是对偶于 $\partial_x$ 的元素。李括号由旧公式给出（见
\S6.2.1）；容易看出这个公式（依赖于 $\nu$）与坐标无关。

注意这个正合列没有自然的、与坐标无关的裂解。但是存在自然裂解
$\mathfrak h\otimes\mathcal O_x\to\widehat{\mathfrak h}_{\nu,x}$（像为
圆盘 $D_x$ 上 $\lambda$-联络空间的正交补）。因此定义
$\widehat{\mathfrak h}_{\nu,x}$-模
\[
    \pi^\nu_{\lambda,x}
    =\Ind_{\mathfrak h\otimes\mathcal O_x\oplus\C\mathbf 1}
    ^{\widehat{\mathfrak h}_{\nu,x}}\lambda.
\]
由上一节得到的 $\operatorname{Aut}\mathcal O$ 在李代数
$\widehat{\mathfrak h}_\nu$ 上作用的描述，立刻推出
\[
    \pi^\nu_{\lambda,x}
    =\mathcal Aut_x\times_{\operatorname{Aut}\mathcal O}\pi^\nu_\lambda.
\]

现在可以陈述自由场同态 $w_\kappa$ 的无坐标版本。

\renewcommand{\theproposition}{6.2.4}
\begin{proposition}
    对任意水平 $\kappa$ 与任意点 $x$，存在
    $\gghat_{\kappa,x}$-模的自然同态
    $V_\kappa(\fg)_x\to M_{\mathfrak g,x}\otimes
    \pi^{\kappa-\kappa_c}_{0,x}$。此外，对任意最高权 $\lambda\in\mathfrak h^*$，
    \[
        W_{\lambda,\kappa,x}=M_{\mathfrak g,x}\otimes
        \pi^{\kappa-\kappa_c}_{0,x}
    \]
    带有自然的 $\gghat_{\kappa,x}$-模结构。
\end{proposition}
% [Tutor Note: 命题中 $W_{\lambda,\kappa,x}$ 的定义里，Fock 模的下标印作 $0$（即 $\pi^{\kappa-\kappa_c}_{0,x}$），但既然参数中有最高权 $\lambda$，按语境应为 $\pi^{\kappa-\kappa_c}_{\lambda,x}$。按原文照录，并在报告中登记为疑似笔误。]

\begin{proof}
    由命题 6.2.2 与推论 6.2.3，把 $w_\kappa$ 看作
    $\gghat_{\kappa,x}$-模的同态时，它与两边的
    $\operatorname{Aut}\mathcal O$ 作用相容。因此可以用 torsor
    $\mathcal Aut_x$ 对它作扭变。按上面的讨论，得到命题的第一个论断。同样，
    用 $\mathcal Aut_x$ 扭变 $\gghat$ 在 $W_{\lambda,\kappa}$ 上的作用，得到
    第二个论断。
\end{proof}

特别地，当 $\nu=0$ 时，李代数 $\widehat{\mathfrak h}_{0,x}$ 是交换的。
$\Omega^{-\rho}$ 在去心圆盘 $D^\times_x$ 上的任何联络 $\nabla$ 定义了
$\widehat{\mathfrak h}_{0,x}$ 上的线性泛函，从而是把
$\widehat{\mathfrak h}_{0,x}$ 看作交换李代数时的一维表示
$\C_\nabla$。若取同构 $\cK_x\simeq\C(\!(t)\!)$，则联络由公式
$\nabla=\partial_t+\chi(t)$ 给出，其中 $\chi(t)\in\mathfrak h^*(\!(t)\!)$。
生成元 $b_{i,n}$ 在 $\C_\nabla$ 上的作用由公式
\[
    b_{i,n}\mapsto\int\langle\chi(z),h_i\rangle z^n dz
\]
给出。

由定理 6.1.6，存在顶点代数同态
$w_{\kappa_c}:V_{\kappa_c}(\fg)\to M_{\mathfrak g}\otimes\pi_0$。由推论
6.2.3，它与两边的 $\operatorname{Aut}\mathcal O$ 作用交换。因此相应的李代数
同态 $\gghat\to U(M_{\mathfrak g}\otimes\pi_0)$ 也与
$\operatorname{Aut}\mathcal O$ 作用交换。于是可以用
$\operatorname{Aut}\mathcal O$-torsor $\mathcal Aut_x$ 扭变这个同态。然后得到
李代数同态
\[
    \gghat_{\kappa_c,x}\to U(M_{\mathfrak g}\otimes\pi_0)_x
    \overset{\mathrm{def}}{=}\mathcal Aut_x
    \times_{\operatorname{Aut}\mathcal O}U(M_{\mathfrak g}\otimes\pi_0).
\]

称 $\Gamma_x\oplus\widehat{\mathfrak h}_{0,x}$-模是\textbf{光滑的}，若其中
任一向量都被李子代数
\[
    (\mathfrak n_+\otimes\mathfrak m^N_x)
    \oplus(\mathfrak n^*_+\otimes\mathfrak m^N_x\Omega_{\mathcal O_x})
    \oplus(\mathfrak h\otimes\mathfrak m^N_x)
\]
湮灭，其中 $\mathfrak m$ 是 $\mathcal O_x$ 的极大理想，$N\in\Z_+$ 充分大。
显然，任何光滑 $\Gamma\oplus\widehat{\mathfrak h}_0$-模自动是
$U(M_{\mathfrak g}\otimes\pi_0)$-模。因此任何光滑
$\Gamma_x\oplus\widehat{\mathfrak h}_{0,x}$-模自动是
$U(M_{\mathfrak g}\otimes\pi_0)_x$-模，从而是
$\gghat_{\kappa_c,x}$-模。

\renewcommand{\theproposition}{6.2.5}
\begin{proposition}
    对去心圆盘 $D^\times_x=\Spec\cK_x$ 上 $^LH$-丛 $\Omega^{-\rho}$ 上的
    任何联络，$M_{\mathfrak g,x}$ 上都有典范的
    $\gghat_{\kappa_c,x}$-模结构。
\end{proposition}

\begin{proof}
    注意 $M_{\mathfrak g,x}$ 是光滑 $\Gamma_x$-模，而 $\C_\nabla$ 对
    $D^\times_x$ 上 $^LH$-丛 $\Omega^{-\rho}$ 的任何联络 $\nabla$ 都是光滑
    $\widehat{\mathfrak h}_{0,x}$-模。取这两个模的张量积，得到
    $\gghat_{\kappa_c,x}$-模，它作为向量空间同构于 $M_{\mathfrak g,x}$。
\end{proof}

于是得到一族由 $^LH$-丛 $\Omega^{-\rho}$ 在去心圆盘 $D^\times_x$ 上的联络
参数化的 $\gghat_{\kappa_c,x}$-模。若取同构
$\cK_x\simeq\C(\!(t)\!)$ 并把联络写成
$\nabla=\partial_t+\chi(t)$，则对应于 $\nabla$ 的模正是推论 6.1.5 中引入的
Wakimoto 模 $W_{\chi(z)}$。

\section{半无限抛物诱导}

\setcounter{theorem}{0}

\begin{keypoints}
    \item \textbf{核心思想：} 把 Borel 子代数换成任意抛物子代数 $\mathfrak p$，
    把 Heisenberg 子代数换成 $\mathfrak p$ 的 Levi 子代数圈代数的中心扩张，
    得到“半无限抛物诱导”函子。
    \item \textbf{数学工具：} Levi 分解 $\mathfrak p=\mathfrak m\oplus\mathfrak r$、
    齐性空间 $G/P$ 与 $G/R$ 的局部化、Weyl 代数 $\cA^{\mathfrak g,\mathfrak p}$、
    二上循环与上同调、Verma 模的特征标。
    \item \textbf{关键结论：} 定理 6.3.1（半无限抛物诱导的存在性）、推论 6.3.2
    （函子性）、命题 6.3.3（$W^+_{0,\kappa_c}\cong\mathbb M_{0,\kappa_c}$）、
    引理 6.3.4（$\widetilde{\mathfrak b}_+$-不变向量）。
    \item \textbf{记号提醒：}$\Delta'_+$ 是不出现在 $\mathfrak p$ 的根空间分解中的
    正根；$\kappa_{i,c}$ 是 $\mathfrak m_i$ 上的临界内积；$\pi^{\kappa_0}_0$ 是
    Fock 表示。
\end{keypoints}

本节推广 Wakimoto 模的构造：用 $\fg$ 的任意抛物子代数代替 Borel 子代数。

\subsection{Wakimoto 模作为诱导表示}

上面给出的 Wakimoto 模构造可以概括如下：对 Heisenberg 李代数
$\widehat{\mathfrak h}_\kappa$ 的每个表示 $N$，我们在
$M_{\mathfrak g}\otimes N$ 上构造了 $\gghat_{\kappa+\kappa_c}$-模结构。
步骤是：先把 $\widehat{\mathfrak h}_\kappa$-模零延拓到
$\hat{\mathfrak b}_{-,\kappa}$，然后施加一个可视为从
$\hat{\mathfrak b}_{-,\kappa}$-模到 $\gghat_{\kappa+\kappa_c}$-模的
半无限类比诱导函子。与通常诱导不同，这个构造的一个重要特征是水平被
$\kappa_c$ 平移。特别地，若从 $\widehat{\mathfrak h}_0$-模（等价地，交换李代数
$L\mathfrak h$ 的表示）出发，则得到临界水平而非水平 $0$ 的
$\gghat_{\kappa_c}$-模。例如，可以把此构造应用于交换李代数 $L\mathfrak h$
的不可约光滑表示。这些表示是一维的，与（拓扑）对偶空间
$(L\mathfrak h)^*\simeq\mathfrak h^*(\!(t)\!)dt$ 中的元素 $\chi(t)$ 一一
对应。结果得到上面引入的临界水平 Wakimoto 模 $W_{\chi(t)}$。考察 $\chi(t)$
在坐标变换群 $\operatorname{Aut}\mathcal O$ 作用下的变换性质，我们发现
$\chi(t)$ 实际上不按一次形式变换，而按某个特定 $^LH$-丛上的联络变换。这个
“反常”是“半无限”构造的典型特征。

相反，若 $\kappa\ne0$，$\widehat{\mathfrak h}_\kappa$ 的不可约光滑模正是 Fock
表示 $\pi^\kappa_\lambda$。对每个这样的表示，我们附上一个
$\gghat_{\kappa+\kappa_c}$-模 $W_{\lambda,\kappa+\kappa_c}$。

现在用任意抛物子代数 $\mathfrak p$ 及其 Levi 商 $\mathfrak m$ 代替 Borel
子代数 $\mathfrak b$ 及其 Levi 商 $\mathfrak h$，来推广这一构造。于是希望把
圈代数 $L\mathfrak m$ 的某个中心扩张上的模附上一个 $\gghat$-模。结果确实可以
做到，只要选取 $L\mathfrak m$ 的适当中心扩张。称所得的 $\gghat$-模为
对应于 $\mathfrak p$ 的\textbf{广义 Wakimoto 模}。于是得到从光滑
$\widehat{\mathfrak m}$-模范畴到光滑 $\gghat$-模范畴的函子。自然地称之为
\textbf{半无限抛物诱导}函子（类比于约化群表示的类似构造）。

\subsection{主要结果}

设 $\mathfrak p$ 是 $\fg$ 的抛物子代数。假定 $\mathfrak p$ 包含下 Borel
子代数 $\mathfrak b_-$（从而特别地 $\mathfrak p$ 包含
$\mathfrak h\subset\mathfrak b_-$）。设
\begin{equation}
    \mathfrak p=\mathfrak m\oplus\mathfrak r
    \tag{6.3-1}
\end{equation}
是 $\mathfrak p$ 的 Levi 分解，其中 $\mathfrak m$ 是包含 $\mathfrak h$ 的
Levi 子代数，$\mathfrak r$ 是 $\mathfrak p$ 的幂零根。进一步设
\[
    \mathfrak m=\bigoplus_{i=1}^s\mathfrak m_i\oplus\mathfrak m_0
\]
是 $\mathfrak m$ 分解为单李子代数 $\mathfrak m_i$（$i=1,\dots,s$）与阿贝尔
子代数 $\mathfrak m_0$ 的直和，且这些直和项关于 $\fg$ 上的内积两两正交。用
$\kappa_{i,c}$ 表示 $\mathfrak m_i$（$i=1,\dots,s$）上的临界内积，其定义与
\S5.3.3 相同。还置 $\kappa_{0,c}=0$。

给定一族 $\mathfrak m_i$（$i=0,\dots,s$）上的不变内积 $\kappa_i$，得到
$\mathfrak m$ 上的不变内积。设 $\widehat{\mathfrak m}_{(\kappa_i)}$ 是相应
的仿射 Kac--Moody 代数，即 $L\mathfrak m$ 的一维中心扩张，其交换关系由公式
(1.3-4) 给出。用 $V_{\kappa_i}(\mathfrak m_i)$（$i=1,\dots,s$）表示
$\widehat{\mathfrak m}_{i,\kappa_i}$ 上的真空模，顶点代数结构如 \S6.1.1
所定义。还用 $V_{\kappa_0}(\mathfrak m_0)$ 表示 Heisenberg 李代数
$\mathfrak m_{\kappa_0}$ 的 Fock 表示 $\pi^{\kappa_0}_0$，其顶点代数结构如
\S6.2.1 所定义。设
\[
    V_{(\kappa_i)}(\mathfrak m)\overset{\mathrm{def}}{=}
    \bigotimes_{i=0}^sV_{\kappa_i}(\mathfrak m_i)
\]
是 $\widehat{\mathfrak m}_{(\kappa_i)}$ 上的真空模，带有张量积顶点代数
结构。

用 $\Delta'_+$ 表示 $\fg$ 中不出现于 $\mathfrak p$ 的根空间分解中的正根集
（注意按假设 $\mathfrak b_-\subset\mathfrak p$，所有负根都出现）。设
$\cA^{\mathfrak g,\mathfrak p}$ 是由 $a_{\alpha,n}$、$a^*_{\alpha,n}$
（$\alpha\in\Delta'_+$，$n\in\Z$）生成、满足关系 (5.3-3) 的 Weyl 代数。设
$M_{\mathfrak g,\mathfrak p}$ 是由向量 $|0\rangle$ 生成的
$\cA^{\mathfrak g,\mathfrak p}$ 的 Fock 表示，满足
\[
    a_{\alpha,n}|0\rangle=0,\quad n\ge0;\qquad
    a^*_{\alpha,n}|0\rangle=0,\quad n>0.
\]
则 $M_{\mathfrak g,\mathfrak p}$ 带有如 \S5.4.1 所定义的顶点代数结构。

我们有定理 6.2.1 的下述类比。

\renewcommand{\thetheorem}{6.3.1}
\begin{theorem}
    设 $\kappa_i$（$i=0,\dots,s$）是一族内积，使得存在 $\fg$ 上的内积
    $\kappa$，它在 $\mathfrak m_i$ 上的限制等于 $\kappa_i-\kappa_{i,c}$
    （对一切 $i=0,\dots,s$）。则存在顶点代数同态
    \[
        w^{\mathfrak p}_\kappa:V_{\kappa+\kappa_c}(\fg)
        \to M_{\mathfrak g,\mathfrak p}\otimes V_{(\kappa_i)}(\mathfrak m).
    \]
\end{theorem}

\begin{proof}
    证明是定理 6.2.1 证明的推广（事实上，当 $\mathfrak p=\mathfrak b_-$ 时
    定理 6.2.1 是定理 6.3.1 的特例）。设 $P$ 是 $G$ 中对应于 $\mathfrak p$
    的李子群，考虑齐性空间 $G/P$。它有开稠密子集
    $\cU_{\mathfrak p}=N_{\mathfrak p}\cdot[1]$，其中 $N_{\mathfrak p}$ 是
    $N_+$ 中对应于子集 $\Delta'_+\subset\Delta_+$ 的子群。用指数映射把
    $\cU_{\mathfrak p}$ 与 $N_{\mathfrak p}$ 及其李代数
    $\mathfrak n_{\mathfrak p}$ 等同。

    置 $L\cU_{\mathfrak p}=\cU_{\mathfrak p}(\!(t)\!)$。与 \S5.3.1 一样
    定义 $L\cU_{\mathfrak p}$ 上的函数与向量场，分别记作
    $\mathrm{Fun}\,L\cU_{\mathfrak p}$ 与
    $\mathrm{Vect}\,L\cU_{\mathfrak p}$。$L\fg$ 在
    $\cU_{\mathfrak p}(\!(t)\!)$ 上的作用给出李代数同态
    $L\fg\to\mathrm{Vect}\,L\cU_{\mathfrak p}$，与前面一样。如下推广这个
    同态。

    考虑商 $G/R$，其中 $R$ 是 $G$ 中对应于 Levi 分解 (6.3-1) 中出现的幂零
    李代数 $\mathfrak r$ 的李子群。有自然投影 $G/R\to G/P$，它是 $M$-丛，
    其中 $M=P/R$ 是 $G$ 中对应于 $\mathfrak m$ 的 Levi 子群。在
    $\cU_{\mathfrak p}\subset G/P$ 上这个丛可以平凡化，因此同构于
    $\cU_{\mathfrak p}\times M$。李代数 $\fg$ 作用在这个丛上，从而作用在
    $\cU_{\mathfrak p}\times M$ 上。结果得到李代数同态
    $\fg\to\mathrm{Vect}(\cU_{\mathfrak p}\times M)$。容易看出它通过同态
    \[
        \fg\to(\mathrm{Vect}\,\cU_{\mathfrak p}\otimes1)
        \oplus(\mathrm{Fun}\,\cU_{\mathfrak p}\otimes\mathfrak m)
    \]
    与 $\mathfrak m\to\mathrm{Vect}\,M$ 分解。

    这个构造的圈版本给出（拓扑）李代数的（连续）同态
    \[
        L\fg\to\mathrm{Vect}\,L\cU_{\mathfrak p}
        \oplus\mathrm{Fun}\,\cU_{\mathfrak p}
        \hat{\otimes}L\mathfrak m.
    \]
    此外，这个同态的像包含在“局部部分”中，即
    $\mathrm{Vect}\,L\cU_{\mathfrak p}$ 的局部部分
    $\cT^{\mathfrak g,\mathfrak p}_{\mathrm{loc}}$（定义见 \S5.4.3）与
    $\mathrm{Fun}\,\cU_{\mathfrak p}\hat{\otimes}L\mathfrak m$ 的局部部分
    $\cI^{\mathfrak g,\mathfrak p}_{\mathrm{loc}}$ 的直和中。按定义，
    $\cI^{\mathfrak g,\mathfrak p}_{\mathrm{loc}}$ 是形式幂级数
    $P(\partial_z^na^*_\alpha(z))J^a(z)$ 的 Fourier 系数张成的空间，其中
    $P$ 是 $a^*_\alpha(z)$（$\alpha\in\Delta'_+$）的微分多项式，
    $J^a\in\mathfrak m$。

    设 $\cA^{\mathfrak g,\mathfrak p}_{0,\mathrm{loc}}$ 与
    $\cA^{\mathfrak g,\mathfrak p}_{\le1,\mathrm{loc}}$ 是 Weyl 代数
    $\cA^{\mathfrak g,\mathfrak p}$ 的局部完备化上自然滤过的第零项与第一项
    （定义见 \S5.3.3）。有李代数的非分裂正合列
    \begin{equation}
        0\to\cA^{\mathfrak g,\mathfrak p}_{0,\mathrm{loc}}
        \to\cA^{\mathfrak g,\mathfrak p}_{\le1,\mathrm{loc}}
% [Tutor Note: “局部部分”与滤过是第 5 章构造的平行版本；正合列 (6.3-2) 一般非分裂，故不能取 $\cT^{\mathfrak g,\mathfrak p}_{\mathrm{loc}}$ 的显然截面。正规序给出向量空间的分裂，而分裂的失败程度正是二上循环 $\omega_{(\kappa_i)}$ 的来源。]
        \to\cT^{\mathfrak g,\mathfrak p}_{\mathrm{loc}}\to0.
        \tag{6.3-2}
    \end{equation}

    置
    \begin{equation}
        \cJ^{\mathfrak g,\mathfrak p}_{\mathrm{loc}}
        \overset{\mathrm{def}}{=}
        \cA^{\mathfrak g,\mathfrak p}_{\le1,\mathrm{loc}}
        \oplus\cI^{\mathfrak g,\mathfrak p}_{\mathrm{loc}},
        \tag{6.3-3}
    \end{equation}
    并注意 $\cJ^{\mathfrak g,\mathfrak p}_{\mathrm{loc}}$ 自然是局部李代数
    $U(M_{\mathfrak g,\mathfrak p}\otimes V_{(\kappa_i)}(\mathfrak m))$ 的
    李子代数。\transnote{卡点：(6.3-2) 为何注明"非分裂"，分裂失败又落到哪里？背景：$\cT^{\mathfrak g,\mathfrak p}_{\mathrm{loc}}$ 是 $\fg$ 的 loop 部分的一阶对象，$\cA^{\mathfrak g,\mathfrak p}_{\le1,\mathrm{loc}}$ 是它的一阶扩张；正规序只给出向量空间截面 $\overline w_{(\kappa_i)}$，而两个一阶元素的交换会产生中心项，故截面不是李代数同态。这个"差"就是二上循环 $\omega_{(\kappa_i)}$，其显式公式见本节 (6.3-5)--(6.3-6)。落点：\S6.3 式 (6.3-2)、(6.3-5)--(6.3-6)。来源：原书 \S6.3 的推导；并见 \S5.4.3、\S5.6.2 的平行构造。}用正规序把正合列 (6.3-2) 作为向量空间分裂，得到线性映射
    $\overline w_{(\kappa_i)}:L\fg\to
    \cJ^{\mathfrak g,\mathfrak p}_{\mathrm{loc}}$。

    需要计算 $\overline w_{(\kappa_i)}$ 不是李代数同态的程度。于是考虑相应的
    线性映射 $\omega_{(\kappa_i)}:\bigwedge^2L\fg\to
    \cJ^{\mathfrak g,\mathfrak p}_{\mathrm{loc}}$，其定义为
    \[
        \omega_{(\kappa_i)}(f,g)
        =\left[\overline w_\kappa(f),\overline w_\kappa(g)\right]
        -\overline w_\kappa([f,g]).
    \]
    与引理 5.6.2 的证明一样显式计算，可知 $\omega_{(\kappa_i)}$ 取值于
    $\cA^{\mathfrak g,\mathfrak p}_{0,\mathrm{loc}}\subset
    \cJ^{\mathfrak g,\mathfrak p}_{\mathrm{loc}}$。此外，
    $\cA^{\mathfrak g,\mathfrak p}_{0,\mathrm{loc}}$ 自然是 $L\fg$-模，
    并且按 $\overline w_{(\kappa_i)}$ 的构造，对任意
    $X\in\cA^{\mathfrak g,\mathfrak p}_{0,\mathrm{loc}}$ 与 $f\in L\fg$ 有
    $[\overline w_{(\kappa_i)}(f),X]=f\cdot X$。这说明
    $\omega_{(\kappa_i)}$ 是以
    $\cA^{\mathfrak g,\mathfrak p}_{0,\mathrm{loc}}$ 为系数的 $L\fg$
    二上循环。按构造它是局部的，即属于
    $C^2_{\mathrm{loc}}(L\fg,\cA^{\mathfrak g,\mathfrak p}_{0,\mathrm{loc}})$
    （定义见 \S5.6.2）。

    沿用引理 5.6.7 证明中的论证，可以证明：若
    $C^2_{\mathrm{loc}}(L\fg,\cA^{\mathfrak g,\mathfrak p}_{0,\mathrm{loc}})$
    中两个上循环在 $\bigwedge^2(L\mathfrak m)$ 上的限制相同，则它们代表
    同一个上同调类。

    下面计算 $\omega_{(\kappa_i)}$ 在 $\bigwedge^2L\mathfrak m$ 上的限制。
    为此在 $L\mathfrak m$ 的元素上计算 $\overline w_{(\kappa_i)}$。设
    $\{J^\alpha\}_{\alpha\in\Delta'_+}$ 是 $\mathfrak n_{\mathfrak p}$ 的
    一组基。李代数 $\mathfrak m$ 在 $\fg$ 上的伴随作用保持
    $\mathfrak n_{\mathfrak p}$，因此得到 $\mathfrak m$ 在
    $\mathfrak n_{\mathfrak p}$ 上的表示 $\rho_{\mathfrak n_{\mathfrak p}}$。
    对任意元素 $A\in\mathfrak m$ 有
    \[
        \rho_{\mathfrak n_{\mathfrak p}}(A)\cdot J^\alpha
        =\sum_{\beta\in\Delta'_+}c^\alpha_\beta(A)J^\beta,
    \]
    其中 $c^\alpha_\beta(A)\in\C$。因此得到下述公式：\transnote{卡点：下面 (6.3-4) 右端只对 $\beta$ 求和，而系数 $c^\alpha_\beta(A)$ 上的 $\alpha$ 是自由指标，这是否漏写？背景：$\rho_{\mathfrak n_{\mathfrak p}}(A)\cdot J^\alpha=\sum_\beta c^\alpha_\beta(A)J^\beta$ 是 $\mathfrak n_{\mathfrak p}$ 的一组基下的矩阵元，上标 $\alpha$ 标记"哪一个基向量的像"；因此 (6.3-4) 应理解为对每个固定的 $\alpha$ 各成立一式（等价于把 $\overline w_{(\kappa_i)}(A(z))$ 的 $\mathfrak n_{\mathfrak p}$-分量逐项写出），或在 $\alpha$ 上再求和并配以对偶基。原书 p.186 的 (6.3-4) 求和范围印作 $\Delta_+$，而 $c^\alpha_\beta$ 只在 $\beta\in\Delta'_+$ 时有定义；这是原书排印问题，译文照录 $\Delta_+$，请读者以定义式为准。落点：\S6.3.1 式 (6.3-4)。来源：原书 p.186 的 $\rho_{\mathfrak n_{\mathfrak p}}$ 定义与其后 (6.3-4)。}
    \begin{equation}
        \overline w_{(\kappa_i)}(A(z))
        =-\sum_{\beta\in\Delta_+}c^\alpha_\beta(A)
        :\!a^*_\beta(z)a_\beta(z)\!:+\widetilde A(z),
        \qquad A\in\mathfrak m,
        \tag{6.3-4}
    \end{equation}
% [Tutor Note: (6.3-4) 的系数印作 $c^\alpha_\beta(A)$，但右端求和只列了 $\beta$；按上文定义 $\rho_{\mathfrak n_{\mathfrak p}}(A)\cdot J^\alpha=\sum_\beta c^\alpha_\beta(A)J^\beta$，这里应对 $\alpha$ 也求和（或固定 $\alpha$ 逐项理解）。此处按原文照录，核对时以定义式为准。]
    其中 $\widetilde A(z)=\sum_{n\in\Z}(A\otimes t^n)z^{-n-1}$，看作
    $\mathfrak m\subset\cJ^{\mathfrak g,\mathfrak p}_{\mathrm{loc}}$ 中
    元素的生成级数。

    设 $\kappa_{\mathfrak n_{\mathfrak p}}$ 是由公式
    \[
        \kappa_{\mathfrak n_{\mathfrak p}}(A,B)
        =\Tr_{\mathfrak n_{\mathfrak p}}\rho_{\mathfrak n_{\mathfrak p}}(A)
        \rho_{\mathfrak n_{\mathfrak p}}(B)
    \]
    定义的 $\mathfrak m$ 上的内积。直接计算级数 (6.3-4) 的系数之间的交换
    关系，得到：对 $A\in\mathfrak m_i$、$B\in\mathfrak m_j$，有
    \begin{equation}
        \omega_{(\kappa_i)}(A_n,B_m)
        =-n\kappa_{\mathfrak n_{\mathfrak p}}(A,B)\delta_{n,-m},
        \tag{6.3-5}
    \end{equation}
    若 $i\ne j$；且
    \begin{equation}
        \omega_{(\kappa_i)}(A_n,B_m)
        =n\left(-\kappa_{\mathfrak n_{\mathfrak p}}(A,B)
        +\kappa_i(A,B)\right)\delta_{n,-m},
        \tag{6.3-6}
    \end{equation}
    若 $i=j$。于是 $\omega_{(\kappa_i)}$ 在
    $\bigwedge^2L\mathfrak m$ 上的限制取值于常数子空间
    $\C\subset\cA^{\mathfrak g,\mathfrak p}_{0,\mathrm{loc}}$。

    设 $\kappa_{\mathfrak g}$ 是 $\fg$ 上的 Killing 型，$\kappa_{\mathfrak m_i}$
    是 $\mathfrak m_i$ 上的 Killing 型（特别地 $\kappa_{\mathfrak m_0}=0$）。
    则有
    \begin{equation}
        \begin{aligned}
            \kappa_{\mathfrak g}(A,B)
            &=\kappa_{\mathfrak m_i}(A,B)
            +2\kappa_{\mathfrak n_{\mathfrak p}}(A,B),
            &\text{若 } i&=j,\\
            \kappa_{\mathfrak g}(A,B)
            &=2\kappa_{\mathfrak n_{\mathfrak p}}(A,B),
            &\text{若 } i&\ne j.
        \end{aligned}
        \tag{6.3-7}
    \end{equation}
    因子 $2$ 来自必须同时计入正根与负根。回忆我们的假设：对一切 $i\ne j$，
    $\mathfrak m_i$ 与 $\mathfrak m_j$ 关于 $\kappa_{\mathfrak g}$ 正交。这
    说明当 $i\ne j$ 时 $\kappa_{\mathfrak n_{\mathfrak p}}(A,B)=0$。

    回忆按定义 $\kappa_c=-\frac12\kappa_{\mathfrak g}$，
    $\kappa_{i,c}=-\frac12\kappa_{\mathfrak m_i}$。因此公式 (6.3-7) 推出
    \[
        \kappa_{\mathfrak n_{\mathfrak p}}|_{\mathfrak m_i}
        =-\kappa_c+\kappa_{i,c}.
    \]
    于是若对某个 $\fg$ 上的不变内积 $\kappa$ 有
    $\kappa_i=\kappa|_{\mathfrak m_i}+\kappa_{i,c}$，则
    \[
        -\kappa_{\mathfrak n_{\mathfrak p}}|_{\mathfrak m_i}+\kappa_i
        =(\kappa+\kappa_c)|_{\mathfrak m_i}.
    \]
    考察公式 (6.3-5) 与 (6.3-6) 便知：若 $\kappa$ 是 $\fg$ 上的不变内积，
    它在 $\mathfrak m_i$ 上的限制对一切 $i=0,\dots,s$ 等于
    $\kappa_i-\kappa_{i,c}$，则二上循环 $\omega_{(\kappa_i)}$ 在
    $\bigwedge^2(L\mathfrak m)$ 上的限制等于 $L\fg$ 上的二上循环
    $\sigma_{\kappa+\kappa_c}$ 在 $\bigwedge^2(L\mathfrak m)$ 上的限制
    （后者代表对应于 $\fg$ 上内积 $\kappa+\kappa_c$ 的一维中心扩张）。

    应用引理 5.6.7 证明中的论证可知：在上述条件下（这些条件正是定理中陈述的
    条件），$L\fg$ 上的二上循环 $\omega_{(\kappa_i)}$ 与二上循环
    $\sigma_{\kappa+\kappa_c}$ 等价。因此，以与定理 5.6.8 证明相同的方式
    得知：在这些条件下，线性映射
    $\overline w_{(\kappa_i)}:L\fg\to
    \cJ^{\mathfrak g,\mathfrak p}_{\mathrm{loc}}$ 可以通过加上
    $C^1_{\mathrm{loc}}(L\fg,\cA^{\mathfrak g,\mathfrak p}_{0,\mathrm{loc}})$
    的一个元素而修改为李代数同态
    \[
        \gghat_{\kappa+\kappa_c}\to\cJ^{\mathfrak g,\mathfrak p}_{\mathrm{loc}}
        \subset U(M_{\mathfrak g,\mathfrak p}\otimes
        V_{(\kappa_i)}(\mathfrak m)).
    \]
    现在引理 6.1.1 说明存在顶点代数同态
    \[
        w^{\mathfrak p}_\kappa:V_{\kappa+\kappa_c}(\fg)
        \to M_{\mathfrak g,\mathfrak p}\otimes V_{(\kappa_i)}(\mathfrak m).
    \]
    证毕。
\end{proof}

称 $\widehat{\mathfrak m}_{(\kappa_i)}$-模是\textbf{光滑的}，若其中任一向量都
被李子代数 $\mathfrak m\otimes t^N\C[[t]]$ 湮灭（$N$ 充分大）。

\renewcommand{\thecorollary}{6.3.2}
\begin{corollary}
    对任意满足定理 6.3.1 条件的光滑
    $\widehat{\mathfrak m}_{(\kappa_i)}$-模 $R$，张量积
    $M_{\mathfrak g,\mathfrak p}\otimes R$ 自然是光滑
    $\gghat_{\kappa+\kappa_c}$-模。存在从光滑
    $\widehat{\mathfrak m}_{(\kappa_i)}$-模范畴到光滑
    $\gghat_{\kappa+\kappa_c}$-模范畴的函子，把模 $R$ 送到
    $M_{\mathfrak g,\mathfrak p}\otimes R$，把
    $\widehat{\mathfrak m}_{(\kappa_i)}$-同态 $R_1\to R_2$ 送到
    $\gghat_{\kappa+\kappa_c}$-同态
    $M_{\mathfrak g,\mathfrak p}\otimes R_1\to
    M_{\mathfrak g,\mathfrak p}\otimes R_2$。
\end{corollary}

称 $\gghat_{\kappa+\kappa_c}$-模 $M_{\mathfrak g,\mathfrak p}\otimes R$ 为
对应于 $R$ 的\textbf{广义 Wakimoto 模}。

考虑特殊情形：$R$ 是 Wakimoto 模 $W_{\lambda_i,\kappa_i}$（在
$\widehat{\mathfrak m}_i$ 上，$i=1,\dots,s$）与 Fock 表示
$\pi^{\kappa_0}_{\lambda_0}$（在 Heisenberg 李代数
$\widehat{\mathfrak m}_0$ 上）的张量积。此时由构造可知，相应的
$\gghat_{\kappa+\kappa_c}$-模
$M_{\mathfrak g,\mathfrak p}\otimes R$ 同构于
$\gghat_{\kappa+\kappa_c}$ 上的 Wakimoto 模
$W_{\lambda,\kappa+\kappa_c}$，其中 $\lambda=(\lambda_i)$。

最后，假设 $\kappa_i$ 的选择使定理 6.3.1 的条件不满足。则定理 6.3.1 证明中
定义的、以 $\cA^{\mathfrak g,\mathfrak p}_{0,\mathrm{loc}}$ 为系数的
$L\fg$ 上的二上循环 $\omega_{(\kappa_i)}$，限制到 $L\mathfrak m$ 上仍给出
以 $\C$ 为系数的 $L\mathfrak m$ 二上循环。但这时它不再等价于 $L\fg$ 上任一
二上循环 $\sigma$ 在 $L\mathfrak m$ 上的限制（后者可由对应于 $\fg$ 上不变
内积 $\nu$ 的上循环 $\sigma_\nu$ 代表）。与引理 5.6.7 证明中相同的论证给出：
以 $\cA^{\mathfrak g,\mathfrak p}_{0,\mathrm{loc}}$ 为系数的 $L\fg$ 上的
二上循环 $\omega_{(\kappa_i)}$ 不可能等价于以
$\C\subset\cA^{\mathfrak g,\mathfrak p}_{0,\mathrm{loc}}$ 为系数的
$L\fg$ 上的二上循环。因此在这种情况下，映射
$\overline w_{(\kappa_i)}:L\fg\to
\cJ^{\mathfrak g,\mathfrak p}_{\mathrm{loc}}$ 不能提升为李代数同态
$\gghat_\nu\to\cJ^{\mathfrak g,\mathfrak p}_{\mathrm{loc}}$（对任何 $\nu$）。
换言之，定理 6.3.1 的条件是这种同态存在的充分必要条件。

\subsection{一般抛物子代数}

迄今为止我们在抛物子代数 $\mathfrak p$ 包含 $\mathfrak b_-$ 的假设下工作。
也可以构造与其他抛物子代数相伴的 Wakimoto 模。下面说明 $\mathfrak p=
\mathfrak b_+$ 情形如何做。

设 $N$ 是顶点代数 $M_{\mathfrak g}\otimes\pi_0^{\kappa-\kappa_c}$ 上的任意
模。$\fg$ 有对合把 $e_i$ 送到 $f_i$、把 $h_i$ 送到 $-h_i$。在此对合下
$\mathfrak b_-$ 变成 $\mathfrak b_+$。这个对合诱导 $\gghat_\kappa$ 上的对合。
于是定理 6.2.1 说明下列公式在 $N$ 上定义 $\gghat_\kappa$-结构（$\mathbf 1$
作用为恒等）：
\[
    \begin{aligned}
        f_i(z)&\mapsto a_{\alpha_i}(z)+\sum_{\beta\in\Delta_+}
        :\!P^i_\beta(a^*_\alpha(z))a_\beta(z)\!:,\\
        h_i(z)&\mapsto\sum_{\beta\in\Delta_+}
        \beta(h_i):\!a^*_\beta(z)a_\beta(z)\!:-b_i(z),\\
        e_i(z)&\mapsto\sum_{\beta\in\Delta_+}
        :\!Q^i_\beta(a^*_\alpha(z))a_\beta(z)\!:\\
        &\qquad+\left(c_i+(\kappa-\kappa_c)(e_i,f_i)\right)
        \partial_za^*_{\alpha_i}(z)+b_i(z)a^*_{\alpha_i}(z),
    \end{aligned}
\]
其中多项式 $P^i_\beta$、$Q^i_\beta$ 由公式 (5.2-12)--(5.2-14) 引入。

为使所得 $\gghat$-模是具有最高权的模，我们如下选取 $N$。设
$M'_{\mathfrak g}$ 是由向量 $|0\rangle'$ 生成的 Weyl 代数
$\cA^{\mathfrak g}$ 的 Fock 表示，满足
\[
    a_{\alpha,n}|0\rangle'=0,\quad n>0;\qquad
    a^*_{\alpha,n}|0\rangle'=0,\quad n\ge0.
\]
取 $N$ 为模 $M'_{\mathfrak g}\otimes\pi^{\kappa-\kappa_c}_{-2\rho-\lambda}$，
其中 $\pi^{\kappa-\kappa_c}_{-2\rho-\lambda}$ 是 \S6.2.1 中定义的
$\pi^{\kappa-\kappa_c}_0$-模。把这个模记作 $W^+_{\lambda,\kappa}$。这是
对应于抛物子代数 $\mathfrak b_+$ 的广义 Wakimoto 模。对向量
\begin{equation}
    |0\rangle'\otimes|-2\rho-\lambda\rangle
    \in M'_{\mathfrak g}\otimes\pi^{\kappa-\kappa_c}_{-2\rho-\lambda}
    =W^+_{\lambda,\kappa}
    \tag{6.3-8}
\end{equation}
我们继续使用记号 $|0\rangle'$。

下述结果把一个特殊的 Wakimoto 模与 Verma 模等同起来，将在 \S8.1.1 中使用。

考虑李代数 $\widehat{\mathfrak n}_+=(\fg\otimes t\C[[t]])\oplus
(\mathfrak n_+\otimes1)$。对 $\lambda\in\mathfrak h^*$，设 $\C_\lambda$ 是
$\widehat{\mathfrak n}_+\oplus(\mathfrak h\otimes1)\oplus\C\mathbf 1$ 的
一维表示，其中 $(\fg\otimes t\C[[t]])\oplus(\mathfrak n_+\otimes1)$ 作用为
$0$，$\mathfrak h\otimes1$ 按 $\lambda$ 作用，$\mathbf 1$ 作用为恒等。定义
水平 $\kappa$、最高权 $\lambda$ 的\textbf{Verma 模} $\mathbb M_{\lambda,\kappa}$
为相应的诱导 $\gghat_\kappa$-模：
\begin{equation}
    \mathbb M_{\lambda,\kappa}
    =\Ind_{\widehat{\mathfrak n}_+\oplus(\mathfrak h\otimes1)\oplus\C\mathbf 1}
    ^{\gghat_\kappa}\C_\lambda.
    \tag{6.3-9}
\end{equation}
张量积中向量 $1\otimes1$ 在 $\mathbb M_{\lambda,\kappa}$ 中的像就是
$\mathbb M_{\lambda,\kappa}$ 的\textbf{最高权向量}，记作
$v_{\lambda,\kappa}$。

\renewcommand{\theproposition}{6.3.3}
\begin{proposition}
    Wakimoto 模 $W^+_{0,\kappa_c}$ 同构于 Verma 模
    $\mathbb M_{0,\kappa_c}$。
\end{proposition}

\begin{proof}
    由公式 (6.3-8) 给出的向量 $|0\rangle'\in W^+_{0,\kappa_c}$ 满足与最高权
    向量 $v_{\lambda,\kappa}\in\mathbb M_{0,\kappa_c}$ 相同的性质：它被
    $\widehat{\mathfrak n}_+$ 湮灭，李代数 $\mathfrak h\otimes1$ 按泛函
    $\lambda$ 作用，$\mathbf 1$ 作用为恒等。因此存在非零同态
    $\mathbb M_{0,\kappa_c}\to W^+_{0,\kappa_c}$，把
    $\mathbb M_{0,\kappa_c}$ 的最高权向量送到 $|0\rangle'\in
    W^+_{0,\kappa_c}$。

    我们先证明模 $\mathbb M_{0,\kappa_c}$ 与 $W^+_{0,\kappa_c}$ 的特征标
    相等。这将把问题归结为证明上述同态是满射。

    回忆 $\gghat_\kappa$-模的特征标概念。设 $\gghat_\kappa$-模 $M$ 带有
    分次算子 $L_0=-t\partial_t$ 的作用，且与它在 $\gghat_\kappa$ 上的作用
    相容。\footnote{注意若 $\kappa\ne\kappa_c$，则任何光滑
    $\gghat_\kappa$-模都通过 Segal--Sugawara 构造带有 Virasoro 代数作用，
    从而特别地带有 $L_0$ 作用。但一般的 $\gghat_{\kappa_c}$-模不必带有
    $L_0$ 作用。} 再设 $L_0$ 与 $\mathfrak h\otimes1\subset\gghat_\kappa$
    在 $M$ 上对角化，且公共特征空间有限维。则定义 $M$ 的特征标为形式级数
    \begin{equation}
        \mathrm{ch}\,M=\sum_{\widehat\lambda\in
        (\mathfrak h\oplus\C L_0)^*}\dim M(\widehat\lambda)
        e^{\widehat\lambda},
        \tag{6.3-10}
    \end{equation}
% [Tutor Note: 特征标的定义要求 $L_0$ 与 $\mathfrak h\otimes1$ 对角化且公共特征空间有限维；临界水平并非所有模都自动带 $L_0$ 作用（见本节脚注），所以命题 6.3.3 的证明要先在 $W^+_{0,\kappa_c}$ 上显式定义 $L_0$。]
    其中 $M(\widehat\lambda)$ 是 $L_0$ 与 $\mathfrak h\otimes1$ 相对于
    $\widehat\lambda:(\mathfrak h\oplus\C L_0)^*\to\C$ 的广义特征空间。\transnote{卡点：特征标的定义为何非要 $L_0$ 与 $\mathfrak h\otimes1$ 同时对角化、且公共特征空间有限维？背景：$\mathrm{ch}\,M$ 是按 $\widehat{\mathfrak h}=(\mathfrak h\otimes1)\oplus\C L_0$ 的权展开的形式级数，只有两者联合对角化时 $\dim M(\widehat\lambda)$ 才有意义；本书把 $\widehat{\mathfrak g}'_\kappa=\C L_0\ltimes\gghat_\kappa$ 的 Cartan 子代数取为 $(\mathfrak h\otimes1)\oplus\C L_0\oplus\C\mathbf 1$，正是为此。临界水平的一般 $\gghat_{\kappa_c}$-模未必带 $L_0$ 作用（见本节脚注），故命题 6.3.3 需先在 $W^+_{0,\kappa_c}$ 上显式定义 $L_0$。落点：\S6.3.2 式 (6.3-10)、\S6.4。来源：原书 p.190；扩展仿射代数的权格约定见 [Kac 1990]。}

    直和 $(\mathfrak h\otimes1)\oplus\C L_0\oplus\C\mathbf 1$ 实际上是
    扩张 Kac--Moody 代数 $\widehat{\mathfrak g}'_\kappa
    =\C L_0\ltimes\gghat_\kappa$ 的 Cartan 子代数（见 [Kac 1990]）。
    这个 Cartan 子代数的对偶空间中的元素称为\textbf{权}。我们考虑中心元
    $\mathbf 1$ 作用为恒等的模中出现的权。因此不妨把这些权看作对偶空间
    $\widetilde{\mathfrak h}=(\mathfrak h\otimes1)\oplus\C L_0$ 的元素，
    从而看作对 $(\lambda,\phi)$，其中 $\lambda\in\mathfrak h^*$，$\phi$ 是
    $-L_0=t\partial_t$ 的值。采用标准记号 $\widehat\delta=(0,1)$。

    $\gghat$ 的正根集自然地是 $\widetilde{\mathfrak h}^*$ 的子集：
    \[
        \widehat\Delta_+
        =\{\alpha+n\delta\mid\alpha\in\Delta_+,n\ge0\}
        \sqcup\{-\alpha+n\delta\mid\alpha\in\Delta_+,n>0\}
        \sqcup\{n\delta\mid n>0\}.
    \]
    前两类根是实根，重数为 $1$。最后一类根是虚根，重数为 $\ell$。

    我们在权集 $\widetilde{\mathfrak h}^*$ 上有自然偏序：
    $\widehat\lambda>\widehat\mu$，若
    $\widehat\lambda-\widehat\mu=\sum_i\widehat\beta_i$，其中
    $\widehat\beta_i$ 是正根。

    设 $\mathbb M_{\lambda,\kappa}$ 是上面定义的 $\gghat_\kappa$ 上的
    Verma 模。把 $\gghat_\kappa$ 的作用唯一地延拓到
    $\widehat{\mathfrak g}'_\kappa$ 上：令
    $L_0\cdot v_{\lambda,\kappa_c}=0$，并利用 $L_0$ 与 $\gghat_\kappa$ 的
    交换关系 $[L_0,A_n]=-nA_n$ 定义 $L_0$ 在
    $\mathbb M_{\lambda,\kappa}$ 其余部分上的作用。所得的模是
    $\widehat{\mathfrak g}'_{\kappa_c}$ 上最高权为
    $\widehat\lambda=(\lambda,0)$ 的 Verma 模，记作
    $\mathbb M_{\widehat\lambda,\kappa}$。

    由 Poincar\'e--Birkhoff--Witt 定理，作为向量空间
    $\mathbb M_{\widehat\lambda,\kappa}$ 同构于 $U(\widehat{\mathfrak n}_-)$，
    其中 $\widehat{\mathfrak n}_-=(\fg\otimes t^{-1}\C[t^{-1}])\oplus
    (\mathfrak n_-\otimes1)$。因此得到 $\mathbb M_{\widehat\lambda,\kappa}$
    的特征标的如下公式：
    \begin{equation}
        \mathrm{ch}\,\mathbb M_{\widehat\lambda,\kappa}
        =e^{\widehat\lambda}\prod_{\widehat\alpha\in\widehat\Delta_+}
        (1-e^{-\widehat\alpha})^{-\mathrm{mult}\,\widehat\alpha},
        \tag{6.3-11}
    \end{equation}
% [Tutor Note: 乘积公式中的 $\mathrm{mult}\,\widehat\alpha$ 取值：实根为 $1$，虚根为 $\ell$。与 (6.3-10) 的 $e^{\widehat\lambda}$ 因子相乘时，注意 $\widehat\lambda$ 同时含 $\mathfrak h^*$ 分量与 $L_0$ 特征值 $\phi$。]
    其中 $\widehat\Delta_+$ 是 $\gghat_\kappa$ 的正根集。

    另一方面，$W^+_{0,\kappa_c}$ 上有来自 \S6.2.4 所述
    $M_{\mathfrak g}\otimes\pi_0$ 上拟共形顶点代数结构的
    $\operatorname{Der}\mathcal O$ 作用，特别地有 $L_0$ 作用。由 \S6.2.5
    得到的公式可知有下列交换关系：
    \[
        [L_0,a_{\alpha,n}]=-na_{\alpha,n},\qquad
        [L_0,a^*_{\alpha,n}]=-na^*_{\alpha,n},\qquad
        [L_0,b_{i,n}]=-nb_{i,n}.
    \]
    这些公式连同要求 $L_0|0\rangle'=0$ 唯一确定 $L_0$ 在
    $W^+_{0,\kappa_c}$ 上的作用。其次，Cartan 子代数
    $\mathfrak h\otimes1\subset\gghat_{\kappa_c}$ 在 $W^+_{0,\kappa_c}$ 上
    有作用，使得
    \[
        [h,a_{\alpha,n}]=\alpha(h)a_{\alpha,n},\qquad
        [h,a^*_{\alpha,n}]=\alpha(h)a^*_{\alpha,n},\qquad
        [h,b_{i,n}]=0
    \]
    对 $h\in\mathfrak h$ 成立。$L_0$ 与 $\mathfrak h\otimes1$ 在
    $|0\rangle\in W^+_{0,\kappa_c}$ 上都作用为 $0$。由于
    $W^+_{0,\kappa_c}$ 的基由单项式
    $a_{\alpha,n}$（$\alpha\in\Delta_+$，$n<0$）、
    $a^*_{\alpha,n}$（$\alpha\in\Delta_+$，$n\le0$）与
    $b_{i,n}$（$i=1,\dots,\ell$，$n<0$）构成，我们发现
    $W^+_{0,\kappa_c}$ 的特征标等于由公式 (6.3-11) 给出的
    $\mathbb M_{0,\kappa_c}$ 的特征标。

    同态 $\mathbb M_{0,\kappa_c}\to W^+_{0,\kappa_c}$ 把两个模上的
    $\widehat{\mathfrak g}'_{\kappa_c}=\C L_0\ltimes\gghat_{\kappa_c}$
    作用交织起来。因此，只要证明同态
    $\mathbb M_{0,\kappa_c}\to W^+_{0,\kappa_c}$ 是满射，或者等价地，
    证明 $W^+_{0,\kappa_c}$ 由公式 (6.3-8) 给出的向量 $|0\rangle'$ 生成，
    命题就得证。

    假设 $W^+_{0,\kappa_c}$ 不由 $|0\rangle'$ 生成。则
    $W^+_{0,\kappa_c}$ 关于李代数
    \[
        \widehat{\mathfrak n}_-=(\mathfrak n_-\otimes1)
        \oplus(\fg\otimes t^{-1}\C[t^{-1}])
    \]
    的余不变量空间的维数大于 $1$，因为除最高权向量的像所张成的一维子空间之外，
    它还必须包含一些向量。这意味着存在 $W^+_{0,\kappa_c}$ 上的齐次线性泛函，
    其权小于最高权 $(0,0)$，且是 $\widehat{\mathfrak n}_-$-不变的。

    于是它特别地关于李子代数
    \[
        L_-\mathfrak b_-=\mathfrak n_-\otimes\C[t^{-1}]
        \oplus\mathfrak h\otimes t^{-1}\C[t^{-1}]
    \]
    不变。因此这个泛函必然通过 $W^+_{0,\kappa_c}$ 关于
    $L_-\mathfrak b_-$ 的余不变量空间分解。

    然而，由 $W^+_{0,\kappa_c}$ 的构造可知，$L_-\mathfrak b_-$ 自由地作用在
    $W^+_{0,\kappa_c}$ 上，且关于这个作用的余不变量空间同构于子空间
    \[
        \C[a^*_{\alpha,n}]_{\alpha\in\Delta_+,n<0}
        \subset W^+_{0,\kappa_c}.
    \]
    事实上，由 $e_\alpha(z)$ 与 $h_i(z)$ 作用的显式公式
    (6.1-2) 与 (6.2-2)（在施加我们的对合后它们变成 $f_\alpha(z)$ 与
    $-h_i(z)$）可知，按字典序排列的单项式
    \[
        \prod_{l_a<0}h_{i_a,l_a}\prod_{m_b\le0}f_{\alpha_b,m_b}
        \prod_{n_c<0}a^*_{\beta_c,n_c}|0\rangle'
    \]
    构成 $W^+_{0,\kappa_c}$ 的基。

    于是得到：$W^+_{0,\kappa_c}$ 上任何 $L_-\mathfrak b_-$-不变泛函完全由它在
    子空间 $\C[a^*_{\alpha,n}]_{\alpha\in\Delta_+,n<0}$ 上的限制确定。因此，
    $W^+_{0,\kappa_c}$ 上权严格小于最高权 $(0,0)$ 的非零
    $L_-\mathfrak b_-$-不变泛函，必然在
    $\C[a^*_{\alpha,n}]_{\alpha\in\Delta_+,n<0}$ 的某个非零权齐次子空间上
    取非零值。但这些子空间的权形如
    \begin{equation}
        -\sum_j(n_j\delta-\beta_j),\qquad n_j>0,\qquad
        \beta_j\in\Delta_+.
        \tag{6.3-12}
    \end{equation}
    由于按假设我们的子空间的权应小于最高权，这个公式中的项数必须非零。

    这说明 $W^+_{0,\kappa_c}$ 必有一个最高权为该形式的不可约商子模。由于
    $W^+_{0,\kappa_c}$ 与 $\mathbb M_{0,\kappa_c}$ 的特征标相等，而不可约
    最高权表示的特征标线性无关（见 [Kac and Kazhdan 1979]），所以
    $\mathbb M_{0,\kappa_c}$ 也有最高权形如 (6.3-12) 的不可约商子模。

    现在回忆 [Kac and Kazhdan 1979] 中描述 Verma 模不可约商子模最高权集的
    Kac--Kazhdan 定理。在所讨论的情形，陈述如下：权
    $\widehat\mu=(\mu,n)$ 作为不可约商子模的最高权出现于
    $\mathbb M_{\widehat\lambda,\kappa}$，其中 $\widehat\lambda=(\lambda,0)$，
    当且仅当 $n\le0$，并且或者 $\mu=\lambda$，或者存在权序列
    $\mu_0,\dots,\mu_m\in\mathfrak h^*$，使得 $\mu_0=\mu$，$\mu_m=0$，
    $\mu_{i+1}=\mu_i\pm m_i\beta_i$（$\beta_i$ 为某些正根，$m_i$ 为正整数），
    且满足
    \begin{equation}
        2(\mu_i+\rho,\beta_i)=m_i(\beta_i,\beta_i)
        \tag{6.3-13}
    \end{equation}
    （这里 $(\cdot,\cdot)$ 是由 $\fg$ 上任一非退化不变内积在
    $\mathfrak h^*$ 上诱导的内积）。

    现在注意方程 (6.3-13) 与分析 $\fg$ 上 Verma 模不可约商子模时出现的
    方程一致，后者的最高权在 外尔群的 $\rho$-平移作用下落在 $\lambda$ 的
    轨道中。换言之，权 $\widehat\mu=(\mu,n)$ 出现在
    $\mathbb M_{0,\kappa_c}$ 的分解中，当且仅当 $n\le0$ 且对 $\fg$ 的
    外尔群的某个元素 $w$ 有 $\mu=w(\rho)-\rho$。但对任何 $w$，权
    $w(\rho)-\rho$ 等于 $\fg$ 的负单根之和。因此 $\mathbb M_{0,\kappa_c}$
    的任何不可约商子模的权都形如
    $-n\delta-\sum_im_i\alpha_i$，$n\ge0$，$m_i\ge0$。这样的权不可能是
    (6.3-12) 的形式。

    因此 $W^+_{0,\kappa_c}$ 由最高权向量生成。于是同态
    $\mathbb M_{0,\kappa_c}\to W^+_{0,\kappa_c}$ 是满射。由于两个模的特征标
    相等，可知 $W^+_{0,\kappa_c}$ 同构于 $\mathbb M_{0,\kappa_c}$。
\end{proof}

在命题 9.5.1 中我们将把这个结果推广到其他最高权的 Verma 模情形。

在 \S8.1.1 中我们还需要一个关于 $W^+_{0,\kappa_c}$ 结构的结果。考虑李代数
\[
    \widetilde{\mathfrak b}_+=(\mathfrak b_+\otimes1)
    \oplus(\fg\otimes t\C[[t]]).
\]

\renewcommand{\thelemma}{6.3.4}
\begin{lemma}
    $W^+_{0,\kappa_c}$ 的 $\widetilde{\mathfrak b}_+$-不变向量空间等于
    $\pi_0\subset W^+_{0,\kappa_c}$。
\end{lemma}

\begin{proof}
    由本节开头给出的 $\gghat_{\kappa_c}$ 在 $W^+_{0,\kappa_c}$ 上的作用
    公式可知，$\pi_0$ 中所有向量都被 $\widetilde{\mathfrak b}_+$ 湮灭。
    下面证明 $W^+_{0,\kappa_c}$ 中没有其他 $\widetilde{\mathfrak b}_+$-不变
    向量。

    $\widetilde{\mathfrak b}_+$-不变向量特别地被李子代数
    $L_+\mathfrak n_-=\mathfrak n_-\otimes t\C[[t]]$ 湮灭。在公式 (6.1-3)
    中我们定义了算子 $e^R_{\alpha,n}$（$\alpha\in\Delta_+$，$n\in\Z$）。
    这些算子生成李代数 $L\mathfrak n_+$ 在 $M_{\mathfrak g}$ 上的右作用，
    它与 $L\mathfrak n_+$ 的（左）作用交换（后者是 $\gghat_{\kappa_c}$ 的
    自由场实现的一部分）。这些算子现在作用在 $W^+_{0,\kappa_c}$ 上，但由于
    我们施加了交换 $\mathfrak n_+$ 与 $\mathfrak n_-$ 的对合，现在把它们
    记作 $f^R_{\alpha,n}$。它们生成李代数 $L\mathfrak n_-(\!(t)\!)$ 的
    作用，该作用与 $L\mathfrak n_-$ 的作用交换（后者是
    $\gghat_{\kappa_c}$ 在 $W^+_{0,\kappa_c}$ 上作用的一部分；见本节开头的
    公式）。

    由这些算子的显式公式容易看出，按字典序排列的单项式
    \[
        \prod_{l_a<0}b_{i_a,l_a}\prod_{m_b\le0}f^R_{\alpha_b,m_b}
        \prod_{n_c<0}a^*_{\alpha_c,n_c}|0\rangle'
    \]
    构成 $W^+_{0,\kappa_c}$ 的基。因此有张量积分解
    \[
        W^+_{0,\kappa_c}
        =\overline W^+_{0,\kappa_c}\otimes W^{+,*}_{0,\kappa_c},
    \]
    其中 $W^{+,*}_{0,\kappa_c}$（相应地，
    $\overline W^+_{0,\kappa_c}$）是仅由 $a^*_{\alpha,n}$（相应地，仅由
    $f^R_{\alpha,m}$ 与 $b_{i,l}$）的单项式张成的空间。此外，由于
    $L_+\mathfrak n_-$ 的作用与 $f^R_{\alpha,m}$ 及 $b_{i,l}$ 交换，我们发现
    $L_+\mathfrak n_-$ 只沿这个张量积分解的第二个因子作用。

    其次，与引理 5.6.4 相同的论证证明：作为 $L_+\mathfrak n_-$-模，
    $W^{+,*}_{0,\kappa_c}$ 同构于有一个生成元的自由模的限制对偶。因此
    $W^{+,*}_{0,\kappa_c}$ 中 $L_+\mathfrak n_-$-不变向量空间是一维的，由
    最高权向量张成。我们断言 $W^+_{0,\kappa_c}$ 的
    $L_+\mathfrak n_-$-不变向量空间等于子空间
    $\overline W^+_{0,\kappa_c}$。

    现在设 $W^+_{0,\kappa_c}$ 中有 $\widetilde{\mathfrak b}_+$-不变向量。
    则它必属于 $\overline W^+_{0,\kappa_c}$。但它还被
    $\widetilde{\mathfrak b}_+$ 的其他元素湮灭，特别地被
    $\mathfrak h\otimes1\subset\widetilde{\mathfrak b}_+$ 湮灭。由于
    \[
        [\mathfrak h\otimes1,a^*_{\alpha,n}]
        =\alpha(h)a^*_{\alpha,n},\qquad h\in\mathfrak h,
    \]
    我们发现 $\overline W^+_{0,\kappa_c}$ 中的向量被 $\mathfrak h\otimes1$
    湮灭当且仅当它属于 $\pi_0$（此时它被整个李子代数
    $\widetilde{\mathfrak b}_+$ 湮灭）。因此 $W^+_{0,\kappa_c}$ 的
    $\widetilde{\mathfrak b}_+$-不变向量空间等于 $\pi_0$。
\end{proof}

\section{附录：Kac--Kazhdan 猜想的证明}

\setcounter{theorem}{0}

\begin{keypoints}
    \item \textbf{核心思想：} 用 Wakimoto 模 $W_{\lambda/t}$ 给出 Kac--Kazhdan
    猜想的一个证明：一般临界权的不可约模特征标等于 Verma 模商去“半无限”
    因子。
    \item \textbf{数学工具：} 扩张仿射 Kac--Moody 代数
    $\widehat{\mathfrak g}'_\kappa=\C L_0\ltimes\gghat_\kappa$、广义 Verma 模、
    Kac--Kazhdan 的 $H^\kappa_{\beta,m}$ 超平面、特征标的线性无关性。
    \item \textbf{关键结论：} 定理 6.4.1：对临界水平的一般权 $\widehat\lambda$，
    $\mathrm{ch}\,L_{\widehat\lambda,\kappa_c}=e^{\widehat\lambda}
    \prod_{\alpha\in\widehat\Delta^{\mathrm{re}}_+}(1-e^{-\alpha})^{-1}$。
    \item \textbf{记号提醒：}$\widehat\delta=(0,1)$，$q=e^{-\delta}$；
    $\widehat\Delta^{\mathrm{re}}_+$ 是实正根集。
\end{keypoints}

作为 Wakimoto 模的应用，我们给出 [Kac and Kazhdan 1979] 中 Kac--Kazhdan
猜想的一个证明，遵循 [Frenkel 1991; 2005b]。

我们将使用扩张仿射 Kac--Moody 代数
$\widehat{\mathfrak g}'_\kappa=\C L_0\ltimes\gghat_\kappa$ 与扩张 Cartan
子代数 $\widehat{\mathfrak h}=\C L_0\oplus(\mathfrak h\otimes1)$ 的权，如
命题 6.3.3 证明中所述。已知 $\gghat_\kappa$ 上的 Verma 模
$\mathbb M_{\widehat\lambda,\kappa}$ 有唯一的不可约商，记作
$L_{\widehat\lambda,\kappa}$。

回忆 [Kac and Kazhdan 1979] 对任意一对 $(\beta,m)$（其中 $\beta$ 是
$\gghat_\kappa$ 的正根，$m$ 是正整数）定义了某个子集
$H^\kappa_{\beta,m}\in\widehat{\mathfrak h}^*$。若 $\beta$ 是实根，则
$H^\kappa_{\beta,m}$ 是 $\widehat{\mathfrak h}^*$ 中的超平面；若 $\beta$ 是
虚根，则 $H^{\kappa_c}_{\beta,m}=\widehat{\mathfrak h}^*$，而当
$\kappa\ne\kappa_c$ 时 $H^\kappa_{\beta,m}=\varnothing$。[Kac and Kazhdan
1979] 证明：$L_{\widehat\lambda,\kappa}$ 是
$\mathbb M_{\widehat\mu,\kappa}$ 的商子模，当且仅当下列条件成立：存在有限
权序列 $\widehat\mu_0,\dots,\widehat\mu_n$，使得
$\widehat\mu_0=\widehat\lambda$，$\widehat\mu_n=\widehat\mu$，
$\widehat\mu_{i+1}=\widehat\mu_i-m_i\beta_i$（$\beta_i$ 为某些正根，$m_i$ 为
正整数），且对一切 $i=1,\dots,n$ 有
$\widehat\mu_i\in H^\kappa_{\beta_i,m_i}$。

用 $\widehat\Delta^{\mathrm{re}}_+$ 表示 $\gghat$ 的正实根集（见命题 6.3.3
的证明）。称权 $\widehat\lambda$ 为临界水平的\textbf{一般权}，若它不属于任何
超平面 $H^\kappa_{\beta,m}$，$\beta\in\widehat\Delta^{\mathrm{re}}_+$。由上述
条件容易看出：$\widehat\lambda$ 是临界水平的一般权，当且仅当
$\mathbb M_{\widehat\lambda,\kappa_c}$ 的不可约商子模的最高权都形如
$\widehat\lambda-n\delta$，其中 $n$ 是非负整数（即它们的 $\mathfrak h^*$
分量等于 $\widehat\lambda$ 的 $\mathfrak h^*$ 分量）。\transnote{卡点：何谓"一般权"，为何只排除实根的超平面 $H^\kappa_{\beta,m}$、$\beta\in\widehat\Delta^{\mathrm{re}}_+$？背景：由 Kac--Kazhdan 的子商判据，权 $\widehat\mu$ 与 $\widehat\lambda$ 只差 $\delta$ 的整数倍时二者的不可约商才可能相同；而临界水平的虚根超平面整体退化为 $\widehat{\mathfrak h}^*$（即对虚根 $\beta$ 没有任何排除效果，见上文 $H^{\kappa_c}_{\beta,m}=\widehat{\mathfrak h}^*$），故"一般"只需避开实根超平面。落点：\S6.4、原书 p.194--195。来源：原书 p.194--195 的 Kac--Kazhdan 判据与 [Kac and Kazhdan 1979]。}

% [Tutor Note: “一般权”是 Kac--Kazhdan 理论中的关键开条件：临界水平的虚根超平面整体退化为 $\widehat{\mathfrak h}^*$，故“一般”只排除实根超平面。定理 6.4.1 的证明正是用 Wakimoto 模把不可约性化归为“无奇异向量 + 由最高权向量生成”两条。]
下述论断就是非扭仿射 Kac--Moody 代数的 Kac--Kazhdan 猜想。

\renewcommand{\thetheorem}{6.4.1}
\begin{theorem}
    对临界水平的一般权 $\widehat\lambda$，
    \[
        \mathrm{ch}\,L_{\widehat\lambda,\kappa_c}
        =e^{\widehat\lambda}\prod_{\alpha\in
        \widehat\Delta^{\mathrm{re}}_+}(1-e^{-\alpha})^{-1}.
    \]
\end{theorem}

\begin{proof}
    不妨假设 $\widehat\lambda=(\lambda,0)$。

    利用 $M_{\mathfrak g}$ 上的顶点代数分次，在 Wakimoto 模
    $W_{\chi(t)}$ 上引入分次算子 $L_0$。由定理 6.1.6 中定义
    $W_{\chi(t)}$ 上 $\gghat_{\kappa_c}$-作用的公式可知，这个作用与分次
    相容，当且仅当 $\chi(t)=\lambda/t$，其中 $\lambda\in\mathfrak h^*$。
    此时
    \[
        \mathrm{ch}\,W_{\lambda/t}
        =e^{\widehat\lambda}\prod_{\alpha\in
        \widehat\Delta^{\mathrm{re}}_+}(1-e^{-\alpha})^{-1},
    \]
    其中 $\widehat\lambda=(\lambda,0)$。因此为证明定理，只需证明：若
    $\widehat\lambda$ 是临界水平的一般权，则 $W_{\lambda/t}$ 不可约。
    假设不然，则或者 $W_{\lambda/t}$ 含有奇异向量（即被李子代数
    $\widehat{\mathfrak n}_+=(\fg\otimes t\C[[t]])\oplus
    (\mathfrak n_+\otimes1)$ 湮灭、且不是最高权向量倍数的向量），或者
    $W_{\lambda/t}$ 不由其最高权向量生成。

    假设 $W_{\lambda/t}$ 含有不是最高权向量倍数的奇异向量。这样的向量必被
    李子代数 $L_+\mathfrak n_+=\mathfrak n_+[[t]]$ 湮灭。我们在注 6.1.1 中
    引入了 $\mathfrak n_+(\!(t)\!)$ 在 $M_{\mathfrak g}$ 上的右作用，它与
    左作用交换。\transnote{卡点："注 6.1.1"在本章找不到——6.1.1 是引理，全章没有编号为 6.1.1 的注。背景：原书 p.195 确实写作 "Remark 6.1.1"，但 $\mathfrak n_+(\!(t)\!)$ 在 $M_{\mathfrak g}$ 上的右作用的引入实际见 \S5.2.5（原书 p.168 亦明写 "see Section 5.2.5"），其显式形式正是本节用到的 (6.1-1)--(6.1-3)。落点：\S6.4 的证明、原书 p.195 与 p.168。来源：原书 p.195 与 p.168 的交叉对照。译文照录原书，未改。}由公式 (6.1-3) 显然可知，单项式
    \begin{equation}
        \prod_{n_a<0}e^R_{\alpha_a,n_a}\prod_{m_b\le0}
        a^*_{\alpha_b,m_b}|0\rangle
        \tag{6.4-1}
    \end{equation}
    构成 $M_{\mathfrak g}$ 的基。以与引理 6.3.4 证明相同的方式可证：
    $W_{0,\kappa_c}$ 的 $L_+\mathfrak n_+$-不变向量空间等于
    $W_{0,\kappa_c}$ 中由不含 $a^*_{\alpha,m}$ 的单项式 (6.4-1) 张成的
    子空间 $\overline W_{0,\kappa_c}$。\transnote{卡点：为什么 $L_+\mathfrak n_+=\mathfrak n_+[[t]]$ 的不变子空间恰好等于 $\overline W_{0,\kappa_c}$？背景：(6.4-1) 的单项式含有 $a^*_{\alpha_b,m_b}$（$m_b\le0$）与 $e^R_{\alpha_a,n_a}$（$n_a<0$）两类因子，而 $L_+\mathfrak n_+$ 由 $\mathfrak n_+[[t]]$ 中非负模式组成，作用于 $a^*_{\alpha,m}$（$m\le0$）即得零，只可能打到 $e^R$ 因子上；与引理 6.3.4 的论证同法可以逐项消去含 $a^*$ 的单项式，故不变子空间正是由不含 $a^*_{\alpha,m}$ 的单项式张成的那一块。落点：\S6.4 的证明、式 (6.4-1)、原书 p.195。来源：原书 p.195 的证明；并见引理 6.3.4 的证明。}

% [Tutor Note: 这一步用到了“$L_+\mathfrak n_+$-不变向量空间等于 $\overline W_{0,\kappa_c}$”的结构性事实（与引理 6.3.4 的证明同法）。奇异向量的权形如 $(\lambda,0)-\sum(n_j\delta-\beta_j)$：$\delta$ 的负号来自 $L_0$ 分次，$\beta_j$ 的正号来自根的方向，两者符号相反是本论证的关键。]
    特别地，我们发现 $W_{\lambda/t}$ 的任何不等于最高权向量的奇异向量的权
    都形如 $(\lambda,0)-\sum_j(n_j\delta-\beta_j)$，其中 $n_j>0$，每个
    $\beta_j$ 是 $\fg$ 的正根。于是 $W_{\lambda/t}$ 含有此权的不可约商
    子模。

    现在注意
    \[
        \mathrm{ch}\,\mathbb M_{(\lambda,0),\kappa_c}
        =\prod_{n>0}(1-q^n)^{-\ell}\cdot
% [Tutor Note: 特征标关系 $\mathrm{ch}\,\mathbb M_{(\lambda,0),\kappa_c}=\prod_{n>0}(1-q^n)^{-\ell}\,\mathrm{ch}\,W_{\lambda/t}$ 反映了半无限构造“丢掉 $\ell$ 族虚根方向”的效应；与命题 6.3.3 的证明对比可见，同一乘积公式在临界水平精确成立。]
        \mathrm{ch}\,W_{\lambda/t},
    \]
    其中 $q=e^{-\delta}$（见命题 6.3.3 的证明）。若不可约模
    $L_{\widehat\mu,\kappa_c}$ 作为 $W_{\lambda/t}$ 的商子模出现，则它出现在
    $\mathrm{ch}\,W_{\lambda/t}$ 分解为不可约表示特征标之和的分解中，从而
    出现在 $\mathrm{ch}\,\mathbb M_{(\lambda,0),\kappa_c}$ 的分解中。由于
    不可约表示的特征标线性无关（见 [Kac and Kazhdan 1979]），这说明
    $L_{\widehat\mu,\kappa_c}$ 是 $\mathbb M_{(\lambda,0),\kappa_c}$ 的
    不可约商子模。但这与假设“$(\lambda,0)$ 是临界水平的一般权”矛盾。因此
    $W_{\lambda/t}$ 不含除最高权向量倍数之外的任何奇异向量。

    其次，假设 $W_{\lambda/t}$ 不由其最高权向量生成。则存在
    $W_{\lambda/t}$ 上的齐次线性泛函，其权小于最高权，且在
    $\widehat{\mathfrak n}_-=(\fg\otimes t^{-1}\C[t^{-1}])\oplus
    (\mathfrak n_-\otimes1)$ 下不变，特别地在其李子代数
    $L_-\mathfrak n_+=\mathfrak n_+\otimes t^{-1}\C[t^{-1}]$ 下不变。因此
    这个泛函通过 $W_{\lambda/t}$ 关于 $L_-\mathfrak n_+$ 的余不变量空间
    分解。但 $L_-\mathfrak n_+$ 自由地作用在 $W_{\lambda/t}$ 上，余不变量
    空间同构于 $W_{\lambda/t}$ 的子空间
    $\C[a^*_{\alpha,n}]_{\alpha\in\Delta_+,n\le0}$。于是这个泛函的权形如
    $(\lambda,0)-\sum_j(n_j\delta+\beta_j)$，其中 $n_j\ge0$，每个
    $\beta_j$ 是 $\fg$ 的正根。与上面一样，这与 $\lambda$ 是一般权的假设
    矛盾。因此 $W_{\lambda/t}$ 由其最高权向量生成。我们又知道它不含除最高权
    向量之外的奇异向量。因此 $W_{\lambda/t}$ 不可约。证毕。
\end{proof}

% \endgroup 恢复 preamble.tex 中定理类环境的默认定义，不影响后续各章。
\endgroup



